\documentclass[11pt, a4paper]{report}

\usepackage[utf8]{inputenc}
\usepackage[T1]{fontenc}
\usepackage{lmodern}
\usepackage{amsmath, amssymb, amsthm}
\usepackage{geometry}
\usepackage{enumerate}
\usepackage{hyperref}
\usepackage{mathtools}
\usepackage{microtype}
\usepackage{array,longtable}

\geometry{margin=1in}
\setlength{\emergencystretch}{3em}

\DeclareUnicodeCharacter{2014}{, }

\title{\textbf{The Thiele Machine}\\ \large A Complete Mathematical Specification\\[6pt] \normalsize Derived from the Fully Machine-Checked Coq Corpus}
\author{Devon Thiele}
\date{July 2026}

% Environments
\theoremstyle{definition}
\newtheorem{definition}{Definition}[chapter]
\newtheorem{axiom}{Axiom}[chapter]
\newtheorem{example}{Example}[chapter]
\newtheorem{remark}{Remark}[chapter]

\theoremstyle{plain}
\newtheorem{theorem}{Theorem}[chapter]
\newtheorem{lemma}[theorem]{Lemma}
\newtheorem{corollary}[theorem]{Corollary}
\newtheorem{proposition}[theorem]{Proposition}

% Macros
\newcommand{\N}{\mathbb{N}}
\newcommand{\Z}{\mathbb{Z}}
\newcommand{\Reals}{\mathbb{R}}
\newcommand{\B}{\mathbb{B}}
\newcommand{\Sspace}{\mathcal{S}}
\newcommand{\Ispace}{\mathcal{I}}
\newcommand{\Gspace}{\mathcal{G}}
\newcommand{\Mspace}{\mathcal{M}}
\newcommand{\Rspace}{\mathcal{R}}
\newcommand{\muled}{\mu_{\text{ledger}}}
\newcommand{\trace}{\tau}
\newcommand{\Tr}{\operatorname{Tr}}

\begin{document}

\maketitle

\noindent\textbf{The Thiele Machine is an argument about what an account of computation should preserve.}

Certification is one concrete witness of a broader investigation of structural state, observation, and cost. The abstract laws, theorems, and executable realizations belong to that investigation. Selected projections lose distinctions that the richer state can express; cost laws constrain specified events under explicit hypotheses.

A conventional encoding can preserve those distinctions. A projection that discards them cannot recover them in general. The stronger proposal that these distinctions belong in a necessary foundation of computation remains a research claim. Neither a physical cost metric nor infinitely many independent geometric dimensions follows from the existence of many observables. The status of each result is stated below.

\begin{abstract}
This is the complete mathematical specification of the Thiele Machine: a formal model of computation where structural information carries an explicit, monotone cost $\mu$ that starts at zero and never decreases. Everything here is built from the machine-checked Coq corpus. You can reconstruct the machine from what's in this document alone.

Coq establishes that conclusions follow from axioms. It does \emph{not} establish that axioms model physical reality. That distinction matters and I keep it explicit throughout. Every result is tagged as one of three types: \textbf{(S)}~structural theorems about the machine as a mathematical object, \textbf{(C)}~conditional derivations where a physics conclusion follows from a stated axiom, and \textbf{(R)}~consistency relations that check internal coherence but are not independent predictions.

The core claims are structural (S): No Free Insight; $\mu$-conservation; $\mu$-hierarchy; $\mu$-ledger necessity (five-way independence classification); categorical separation; strict Turing subsumption; the classical CHSH bound ($|S| \leq 2$); the algebraic Tsirelson bound ($|S| \leq 2\sqrt{2}$); and the biconditional with no project-local axioms \texttt{column\_\allowbreak{}contractive\_\allowbreak{}iff\_\allowbreak{}quantum\_\allowbreak{}realizable}. The CHSH algebraic result is pure math. The Landauer connection is motivation. Geometry: only the discrete angle-defect identity for planar triangulations.

The document also covers, with explicit epistemological tags, conditional derivations under named physical axioms (Born rule under linearity, no-cloning, unitarity) and reference explorations (4D simplicial-complex differential geometry, self-reference tower). These are marked (C) or (R) and do not form part of the primary claim ledger. Two named trust boundaries discussed here are: \texttt{bsc\_\allowbreak{}kami\_\allowbreak{}compilation\_\allowbreak{}trusted} (Layer 2 Verilog semantics gap) and \texttt{A\_\allowbreak{}QM} (physical quantum axiom).
\end{abstract}

\tableofcontents

%=============================================================================
% PART I: FOUNDATIONS
%=============================================================================

\chapter{Introduction}
The Thiele Machine makes the cost of structure explicit. Ordinary step-counting assigns one time unit per transition; it does not specify this structural-cost schedule. The Thiele Machine assigns a specific cost $\mu$ to every operation that extracts or asserts structural information. That's the design. Here is where it goes.

The central claim, proven formally: \emph{specified certification and structural-entitlement events have the cost floors proved under their stated hypotheses}. Combined with explicit physical axioms (linearity, information conservation), that constraint produces structural results that parallel quantum mechanics and thermodynamics. Each result's status, whether a structural theorem~\textbf{(S)}, a conditional derivation~\textbf{(C)}, or a consistency relation~\textbf{(R)}, is marked directly. Nothing is hidden.

The document is organized as:
\begin{enumerate}
    \item \textbf{Foundations}: State space, instruction set, operational semantics.
    \item \textbf{Core Theorems}: No Free Insight, $\mu$-ledger necessity, $\mu$-Chaitin, $\mu$-conservation.
    \item \textbf{Computation}: Strict Turing subsumption, oracle pricing as instantiation of universal NFI, $\mu$-complexity classes, $\mu$-cost hierarchy.
    \item \textbf{Symmetry \& Conservation}: Gauge action, Noether's theorem, receipts.
    \item \textbf{Quantum / CHSH algebra}: Tsirelson bound (algebraic), CHSH classical bound, biconditional column-contractive $\iff$ NPA-realizable, three-tier observation taxonomy. Pure polynomial arithmetic, zero physics axioms.
    \item \textbf{Conditional physics}: Born rule (under linearity), unitarity, no-cloning, purification, Landauer (with second-law assumed as a record field), affine EFE family. Each is conditional on stated axioms in its own file.
    \item \textbf{Discrete Geometry}: Discrete angle-defect identity for planar triangulations (not classical Gauss-Bonnet); 4D simplicial complex and Einstein-tensor definitions (reference exploration only, vacuum case is a $0=0$ consistency check).
    \item \textbf{Meta-Theory}: Three-layer Coq/Python/RTL observable contract, categorical structure.
    \item \textbf{Predictions}: Falsifiable predictions tied to specific Coq-verifiable claims.
\end{enumerate}

\section{Epistemological Framework}

Formal verification (Coq) establishes that conclusions follow logically from axioms. It does \textbf{not} establish that axioms model physical reality. This distinction is critical for evaluating the physics claims in Parts~V--VIII.

Every result in this specification falls into one of three categories:

\begin{description}
    \item[\textbf{(S) Structural}] Theorems about the Thiele Machine as a mathematical object. Their conclusions follow from the definitions, explicit hypotheses, and reported library assumptions. \emph{Examples}: $\mu$-conservation, gauge invariance, confluence, halting undecidability, strict Turing subsumption.

    \item[\textbf{(C) Conditional}] Theorems of the form ``if axiom $X$ holds, then physics result $Y$ follows.'' The logical inference is verified; whether axiom $X$ holds in nature is an empirical question. \emph{Examples}: Born rule uniqueness (conditional on linearity), complex necessity (conditional on 2D amplitudes), unitarity (conditional on information conservation).

    \item[\textbf{(R) Consistency}] Algebraic identities or definitional equivalences that verify internal coherence but do not constitute independent predictions. \emph{Example}: the Tsirelson cost definition.
\end{description}

Each major theorem in Parts~V--VIII is annotated with its category. A theorem marked \textbf{(C)} or \textbf{(R)} is not diminished. Conditional derivations and consistency checks are valuable scientific tools. But the reader should not mistake them for derivations from first principles without stated assumptions. See Appendix~B for the complete classification table.

\chapter{Notation and Conventions}
\label{ch:notation}

This chapter exists because several symbols in this document are overloaded,
and at least one ($S$) is used well before it is defined. Read it once; it
resolves ambiguities that would otherwise require guessing from context.

\section{Overloaded symbols}

\begin{description}
  \item[$S$, four distinct meanings.] Disambiguated by context and, where
    ambiguity would be real, by explicit qualification:
    \begin{enumerate}
      \item $\Sspace$ (calligraphic): the \emph{state space}.
      \item $S$ (roman, in operational-semantics chapters): a \emph{particular
        VM state}, an element of $\Sspace$.
      \item $S(n)$ (applied to a natural number): the \emph{successor
        function} on $\N$, i.e.\ $S(n) = n+1$. This is Coq's \texttt{S}
        constructor, and it is the reason cost expressions such as
        $\mathit{flen}\times 8 + S(\mathit{cost})$ are bounded below by $1$:
        $S(n) \geq 1$ for every $n \in \N$. \textbf{This is the notation used
        from Chapter~\ref{ch:isa} onward, before its first formal
        appearance in the cost tables.}
      \item $S$ (in quantum chapters): the \emph{CHSH statistic}, defined in
        Section~\ref{sec:chsh_statistic}. Real-valued.
    \end{enumerate}
  \item[$G$, three distinct meanings.] All three occur, and two occur inside
    the same chapter on emergent geometry:
    \begin{enumerate}
      \item $G$: the \emph{partition graph} component of a VM state.
      \item $G$: \emph{Newton's constant}, fixed to $1/(8\pi)$ as a unit choice.
      \item $G_{\mu\nu}$: the \emph{Einstein tensor}.
    \end{enumerate}
    Where the geometry chapter needs the partition graph it is written
    $G_{\mathrm{part}}$.
  \item[$\mu$, ledger and index, and two different number systems.]
    \begin{enumerate}
      \item $\mu$: the \emph{cost ledger}, a component of VM state. In the
        machine, \textbf{$\mu$ ranges over $\N$}: it is a natural number, it is
        monotone non-decreasing under the step relation, and A2 gives it an
        integer floor.
      \item $\mu, \nu$ as \emph{tensor indices} in the geometry chapter
        ($g_{\mu\nu}$, $T_{\mu\nu}$). Unrelated to the ledger.
      \item $\mu$ in the \textbf{quantum chapters} (measurement cost, the
        no-cloning inequality, the $\mu/I$ ratio) is \textbf{real-valued}.
        This is a genuinely different object from the machine's $\mu$: there
        is no theorem in this development identifying them, and no coercion
        between them is claimed. Where the distinction matters the real-valued
        quantity is written $\mu_{\mathbb{R}}$.
    \end{enumerate}
\end{description}

\section{Terms used before definition}

\begin{description}
  \item[\texttt{CERTIFY}.] The instruction family whose step sets the
    certification bit. Formally $\texttt{CERTIFY}(\Delta\mu)$ has step
    semantics $\mathrm{cert} \mapsto \mathrm{true}$ with cost
    $S(\Delta\mu) \geq 1$; see Section~\ref{sec:certify_def}. Chapters~5,~7
    and~8 depend on it.
  \item[\texttt{irreversible\_\allowbreak{}bits}.] The count of bits whose prior value
    cannot be recovered from the post-state of a step. It is the quantity
    Landauer's principle prices. Used in the thermodynamic corollaries.
  \item[\texttt{HasStructureAddition}.] In this document, read it as the Coq
    definition rather than the looser prose gloss ``a cert-setting or
    equivalent operation'': the certification CSR transitions from zero to
    non-zero along the trace. ``Or equivalent operation'' is not defined and
    should not be relied on.
\end{description}

\chapter{The State Space}

\section{Primitives}
Let $\N = \{0,1,2,\dots\}$. Let $\B = \{0,1\}$. Let $\Sigma$ be a finite alphabet.

\section{Regions and Modules}
\begin{definition}[Region]
A \textbf{Region} $R$ is a finite subset of $\N$, canonically represented as a deduplicated list preserving first-occurrence order: $\mathrm{norm}(L) = \mathrm{nodup}(L)$ removes duplicate elements while preserving the order of their first appearance. (The Python VM applies a stricter canonical form via \texttt{sorted(set(region))}; the Coq specification uses \texttt{nodup Nat.eq\_\allowbreak{}dec} which does not sort.)
\end{definition}

\begin{definition}[Module]
A \textbf{Module} $M = (R_M, A_M)$ consists of a normalized region $R_M \subset \N$ and a set of axioms $A_M \subset \Sigma^*$.
\end{definition}

\section{The Partition Graph}
\begin{definition}[Partition Graph]
A partition graph $G = (N_{id}, \mathcal{T})$ where $N_{id} \in \N$ is the next available module ID and $\mathcal{T}: \{0,\dots,N_{id}-1\} \rightharpoonup \Mspace$ is a finite partial function from IDs to modules.
\end{definition}

\begin{axiom}[Well-Formedness]
$G$ is well-formed iff $\mathrm{dom}(\mathcal{T}) \subseteq \{0,\dots,N_{id}-1\}$.
\end{axiom}

\section{Machine State}
\begin{definition}[VM State]
The complete state is a 12-tuple:
\[ S = (G, C, R, \mathit{Mem}, \mathit{PC}, \mu, T_\mu, \mathit{err}, \ell, m, w, \mathit{cert}) \]
where $G \in \Gspace$ is the partition graph, $C \in \N^4$ comprises the CSRs (certificate address, status, error code, heap base), $R \in \N^{\texttt{REG\_\allowbreak{}COUNT}}$ is the register file, $\mathit{Mem} \in \N^{\texttt{MEM\_\allowbreak{}SIZE}}$ is main memory (the kernel constants \texttt{REG\_\allowbreak{}COUNT} and \texttt{MEM\_\allowbreak{}SIZE} are parametric; the synthesized RTL fixes \texttt{REG\_\allowbreak{}COUNT}~$=16$ and \texttt{MEM\_\allowbreak{}SIZE}~$=128$), $\mathit{PC} \in \N$ is the program counter, $\mu \in \N$ is the $\mu$-ledger, $T_\mu \in \N^{16}$ is the flattened $4\times 4$ $\mu$-tensor (row-major), $\mathit{err} \in \B$ is the global error flag (latching), $\ell \in \N$ is the logic-engine accumulator, $m \in \N$ is the machine status register ($0$~=~Turing mode, $1$~=~Thiele mode), $w \in \N^8$ is the WitnessCount record (\texttt{vm\_\allowbreak{}witness}): eight natural-number counters \texttt{wc\_\allowbreak{}same\_\allowbreak{}ab} and \texttt{wc\_\allowbreak{}diff\_\allowbreak{}ab} for each setting pair $(a,b) \in \{00,01,10,11\}$, and $\mathit{cert} \in \B$ is the certified flag (\texttt{vm\_\allowbreak{}certified}).
\end{definition}

\chapter{Instruction Set and Cost Model}
\label{ch:isa}

\section{The Axiom Scheme A1--A6}
\label{sec:axiom_scheme}

``A2'' is referred to throughout this document, and this is the scheme it
belongs to. These are \emph{conditions on Coq record types}, not
\texttt{Axiom} declarations. The assumption receipt distinguishes project-local declarations from standard-library axioms and explicit theorem hypotheses. Supplying an instance of the record is the only way to
invoke the corresponding theorems, and the 51-opcode VM supplies them by
construction.

\subsection{The consolidated \texttt{CertificationSystem} (A1--A2)}

A \textbf{CertificationSystem} is a tuple
$(\Sigma, \Ispace, \mathrm{step}, \mathrm{cost}, \mathrm{cert})$ where
$\Sigma$ is any state type, $\Ispace$ any instruction type,
$\mathrm{step} : \Sigma \times \Ispace \to \Sigma$,
$\mathrm{cost} : \Ispace \to \N$, and
$\mathrm{cert} : \Sigma \to \mathbb{B}$, subject to:

\begin{description}
  \item[\textbf{A1} (certification is observable).] A total function
    $\mathrm{cert} : \Sigma \to \mathbb{B}$ exists. This holds \emph{by record
    existence}: providing the field is providing the axiom.
  \item[\textbf{A2} (certification is priced; \emph{the load-bearing one}).]
    \begin{equation}
      \forall s \in \Sigma,\ \forall i \in \Ispace,\quad
      \mathrm{cert}(s) = \mathrm{false} \ \wedge\
      \mathrm{cert}(\mathrm{step}(s,i)) = \mathrm{true}
      \ \implies\ \mathrm{cost}(i) \geq 1 .
    \end{equation}
    In Coq this is the field \texttt{cs\_\allowbreak{}cert\_\allowbreak{}costs}. Everything called ``No
    Free Insight'' in this document is A2 lifted from steps to traces by
    induction; the trace-level statement is genuinely proved, but its content
    is A2 plus ``a sum of naturals containing a summand ${\geq}1$ is
    ${\geq}1$''.
\end{description}

\begin{remark}[A2 is a stipulation, and that is the intended reading]
\label{rem:a2_stipulated}
A2 is a condition a substrate may or may not satisfy, not a discovered fact
about computation. Nothing in this document derives it; Section~\ref{sec:f1}
derives it \emph{conditionally} from a Landauer premise plus a named
calibration bridge, which relocates the stipulation rather than removing it.
Claims of the form ``the machine proves certification must cost'' should be
read as ``on any substrate satisfying A2, certification costs''. That is
what \texttt{universal\_\allowbreak{}nfi\_\allowbreak{}any\_\allowbreak{}substrate} states, over an arbitrary state
type and an arbitrary instruction type.
\end{remark}

\subsection{The \texttt{QuantitativeCertificationSystem} extension (A3--A6)}

The quantitative extension lifts the floor from $\geq 1$ to $\geq K$ for an
arbitrary threshold. It adds a witness function
$W : \Sigma \to \N$ and a threshold $K \in \N$, subject to:

\begin{description}
  \item[\textbf{A3} (cost bounds witness growth).]
    $\forall s, i:\ W(s) + \mathrm{cost}(i) \geq W(\mathrm{step}(s,i))$.
    Evidence cannot accumulate faster than it is paid for. Coq:
    \texttt{qcs\_\allowbreak{}cost\_\allowbreak{}bounds\_\allowbreak{}witness}.
  \item[\textbf{A4} (certification requires the threshold).]
    $\forall s:\ \mathrm{cert}(s) = \mathrm{true} \implies W(s) \geq K$. Coq:
    \texttt{qcs\_\allowbreak{}cert\_\allowbreak{}threshold\_\allowbreak{}witness}.
  \item[\textbf{A5} (clean start).] The initial state is uncertified,
    $\mathrm{cert}(s_0) = \mathrm{false}$, and carries no witness,
    $W(s_0) = 0$. Supplied as a hypothesis at use sites rather than as a
    record field.
  \item[\textbf{A6} (threshold is the claim's price).] $K$ is the parameter of
    the conclusion: a run from $s_0$ to a certified state has total cost
    $\geq K$. Not an independent assumption; it names the quantity A3 and
    A4 telescope into.
\end{description}

\begin{remark}[Two numbering systems, and which one this document uses]
\label{rem:two_numberings}
The kernel carries a second, older numbering attached to the
\texttt{AbstractCertMachine} record, in which the cost-on-certification rule is
called \textbf{A4} and \textbf{A2} names ``cert starts unset''. The two systems
agree on content and disagree on labels. \textbf{Throughout this document,
``A2'' means the cost-on-certification rule in the consolidated
\texttt{CertificationSystem} sense stated above.} When reading kernel sources
that mention \texttt{AbstractCertMachine}, translate accordingly.
\end{remark}

\section{Instruction Categories}
The instruction set $\Ispace$ is organized into six families. Every instruction $i$ carries an explicit cost parameter $\Delta\mu_i \in \N$.

\subsection{Structural Operations}
These modify the partition graph $G$:
\begin{itemize}
    \item $\texttt{PNEW}(R, \Delta\mu)$: Create module with region $R$.
    \item $\texttt{PSPLIT}(m, L, R, \Delta\mu)$: Split module $m$ into disjoint regions.
    \item $\texttt{PMERGE}(m_1, m_2, \Delta\mu)$: Merge two modules.
    \item $\texttt{PDISCOVER}(m, E, \Delta\mu)$: Discovery accounting for module $m$. The abstract specification attaches evidence $E$ (list of axioms) to the module; the current hardware-aligned step advances PC and charges $\Delta\mu$ without mutating the graph.
    \item $\texttt{MDLACC}(m, \Delta\mu)$: Minimum-description-length accumulate on module $m$.
\end{itemize}

\subsection{Logical Operations}
These interact with information content:
\begin{itemize}
    \item $\texttt{LASSERT}(\mathit{freg}, \mathit{creg}, \mathit{kind}, \mathit{flen}, \mathit{cost})$: Assert a formula read from \texttt{vm\_\allowbreak{}mem}[\texttt{R[freg]}] (encoded-unit count $\mathit{flen}$, eight concrete bits per unit) against a witness package at \texttt{vm\_\allowbreak{}mem}[\texttt{R[creg]}]. In SAT mode the package contains one satisfying assignment and one falsifying assignment, stored back-to-back; $\Delta\mu = \mathit{flen} \times 8 + S(\mathit{cost}) \geq 1$.
    \item $\texttt{LJOIN}(\mathit{c1reg}, \mathit{c2reg}, \Delta\mu)$: Join certificates from two memory-resident strings addressed by registers.
    \item $\texttt{REVEAL}(m, n, \mathit{cert}, \Delta\mu)$: Reveal $n$ bits of structure from $m$; cost is $n + S(\Delta\mu)$.
    \item $\texttt{EMIT}(m, p, \Delta\mu)$: Emit payload $p$ from module $m$; cost is $|p|_{\mathrm{bits}} + S(\Delta\mu)$.
\end{itemize}

\subsection{Computational Operations (Reversible ALU)}
Reversible register/memory operations:
\begin{itemize}
    \item $\texttt{XFER}(\mathit{dst}, \mathit{src}, \Delta\mu)$: Copy register $\mathit{src}$ to $\mathit{dst}$.
    \item $\texttt{XOR\_\allowbreak{}LOAD}(\mathit{dst}, \mathit{addr}, \Delta\mu)$: Load memory word at $\mathit{addr}$.
    \item $\texttt{XOR\_\allowbreak{}ADD}(\mathit{dst}, \mathit{src}, \Delta\mu)$: $\mathit{dst} \gets \mathit{dst} \oplus \mathit{src}$.
    \item $\texttt{XOR\_\allowbreak{}SWAP}(a, b, \Delta\mu)$: Swap registers $a$ and $b$.
    \item $\texttt{XOR\_\allowbreak{}RANK}(\mathit{dst}, \mathit{src}, \Delta\mu)$: $\mathit{dst} \gets \mathrm{popcount}(\mathit{src})$.
\end{itemize}

\subsection{Standard ALU and Control Flow}
\begin{itemize}
    \item $\texttt{LOAD\_\allowbreak{}IMM}(\mathit{dst}, \mathit{imm}, \Delta\mu)$: Load 32-bit immediate into register.
    \item $\texttt{LOAD}(\mathit{dst}, \mathit{addr}, \Delta\mu)$: Load from main memory.
    \item $\texttt{STORE}(\mathit{addr}, \mathit{src}, \Delta\mu)$: Store to main memory.
    \item $\texttt{ADD}(\mathit{dst}, r_1, r_2, \Delta\mu)$: $\mathit{dst} \gets r_1 + r_2 \pmod{2^{64}}$.
    \item $\texttt{SUB}(\mathit{dst}, r_1, r_2, \Delta\mu)$: $\mathit{dst} \gets r_1 - r_2 \pmod{2^{64}}$.
    \item $\texttt{JUMP}(\mathit{target}, \Delta\mu)$: Unconditional jump.
    \item $\texttt{JNEZ}(r, \mathit{target}, \Delta\mu)$: Jump if $R[r] \neq 0$.
    \item $\texttt{CALL}(\mathit{target}, \Delta\mu)$: Push return address, jump to target.
    \item $\texttt{RET}(\Delta\mu)$: Pop return address, jump back.
\end{itemize}

\subsection{Auxiliary Operations}
\begin{itemize}
    \item $\texttt{CHSH\_\allowbreak{}TRIAL}(x, y, a, b, \Delta\mu)$: Run a CHSH correlation trial; $x,y,a,b \in \{0,1\}$ required.
    \item $\texttt{CHECKPOINT}(\mathit{label}, \Delta\mu)$: Execution checkpoint (charges $\mu$).
    \item $\texttt{READ\_\allowbreak{}PORT}(\mathit{dst}, \mathit{ch}, \mathit{val}, \mathit{bits}, \Delta\mu)$: Read from I/O channel (cert-setter).
    \item $\texttt{WRITE\_\allowbreak{}PORT}(\mathit{ch}, \mathit{src}, \Delta\mu)$: Write to I/O channel.
    \item $\texttt{HEAP\_\allowbreak{}LOAD}(\mathit{dst}, \mathit{addr}, \Delta\mu)$: Load from heap-relative address.
    \item $\texttt{HEAP\_\allowbreak{}STORE}(\mathit{addr}, \mathit{src}, \Delta\mu)$: Store to heap-relative address.
    \item $\texttt{HALT}(\Delta\mu)$: Halt execution.
\end{itemize}

\subsection{Categorical Morphism Operations}
These 7 opcodes (0x27--0x2D) implement the categorical layer over partition modules. See Chapter~\ref{ch:categorical} for laws and proofs.
\begin{itemize}
    \item $\texttt{MORPH}(\mathit{dst}, \mathit{src}, \mathit{tgt}, \mathit{idx}, \Delta\mu)$: Create morphism $\mathit{src} \to \mathit{tgt}$; write morphism ID to $\texttt{regs}[\mathit{dst}]$.
    \item $\texttt{COMPOSE}(\mathit{dst}, f, g, \Delta\mu)$: Compose $f;g$ when $f.\mathrm{tgt} = g.\mathrm{src}$; errors on type mismatch.
    \item $\texttt{MORPH\_\allowbreak{}ID}(\mathit{dst}, m, \Delta\mu)$: Create identity morphism $\mathrm{id}_m$ (diagonal relation).
    \item $\texttt{MORPH\_\allowbreak{}DELETE}(f, \Delta\mu)$: Remove morphism $f$; cascades when modules are split/merged.
    \item $\texttt{MORPH\_\allowbreak{}ASSERT}(f, \mathit{prop}, \mathit{cert}, \Delta\mu)$: \textbf{Cert-setter.} Charges $S(\Delta\mu) \geq 1$ unconditionally. Even failed attempts (morphism does not exist) incur cost $1$. This is the strongest form of the NoFI law: the machine refuses to let even a doomed attempt at certifying a relational claim slip through for free.
    \item $\texttt{MORPH\_\allowbreak{}TENSOR}(\mathit{dst}, f, g, \Delta\mu)$: Tensor product $f \otimes g$ (requires union modules to exist).
    \item $\texttt{MORPH\_\allowbreak{}GET}(\mathit{dst}, f, \mathit{sel}, \Delta\mu)$: Read morphism metadata: selector $0=\mathrm{src}$, $1=\mathrm{tgt}$, $2=\mathrm{coupling\ length}$, $3=\mathrm{is\_\allowbreak{}identity}$.
\end{itemize}

\emph{A reminder before the machinery: this ISA is scaffolding. It is one way to build the substrate, not the substrate itself. Burn all 51 opcodes and the machine is exactly as real, because the machine is the law, not the instruction set.}

The active implementation exposes a 51-opcode ISA: 14 arithmetic/data, 8 memory and control flow, 7 I/O and tensor, 4 partition operations, 3 logic and certification, 8 witness and certification (\texttt{CHSH\_\allowbreak{}TRIAL}, \texttt{CERTIFY}, \texttt{REVEAL}, \texttt{CHSH\_\allowbreak{}LASSERT}, and the four $Q_{1+AB}$ check variants), 7 categorical morphism operations ($14+8+7+4+3+8+7=51$). 47 of the 51 are synth-realized; the four $Q_{1+AB}$ variants live in the Kami HW abstraction with kernel-equivalence proven but are excluded from the synthesized Verilog by silicon budget. This condensed mathematical specification presents the core instruction families relevant to the kernel semantics and cost model. Every instruction carries an explicit $\Delta\mu$ parameter; the function $\texttt{instruction\_\allowbreak{}cost}$ extracts it.

\section{Cost Model}
Every instruction $i$ carries an explicit cost parameter $\Delta\mu_i \in \N$. The function $\texttt{instruction\_\allowbreak{}cost}(i)$ extracts this parameter. The cost is applied via $\texttt{apply\_\allowbreak{}cost}(s, i) = s.\mu + \texttt{instruction\_\allowbreak{}cost}(i)$. Since $\Delta\mu_i : \N$, the $\mu$-ledger is monotonically non-decreasing by construction.

The mandatory positive-cost class is: $\texttt{CERTIFY}$, $\texttt{MORPH\_\allowbreak{}ASSERT}$, and $\texttt{LJOIN}$ (each charges $S(\Delta\mu) \geq 1$); $\texttt{EMIT}$ ($|\mathit{payload}|_{\mathrm{bits}} + S(\Delta\mu)$); $\texttt{REVEAL}$ and $\texttt{READ\_\allowbreak{}PORT}$ ($\mathit{bits} + S(\Delta\mu)$); and $\texttt{LASSERT}$ ($\mathit{flen} \times 8 + S(\Delta\mu)$). These cannot be zero-cost even with $\Delta\mu = 0$. The five \texttt{CHSH\_\allowbreak{}LASSERT} family instructions also charge $S(\Delta\mu)$. Other instructions charge their encoded $\Delta\mu$, which may be zero.

\begin{theorem}[Least joint unit floor, \textbf{(S)}]
For a finite family of event predicates on instructions, the least natural-valued schedule satisfying a unit floor on each event is the indicator that at least one event occurs. This is \texttt{joint\_\allowbreak{}floor\_\allowbreak{}is\_\allowbreak{}least}. Independent Boolean coordinates can change together for one unit under these rules; additivity requires an additional condition.
\end{theorem}

\section{The \texttt{CERTIFY} instruction}
\label{sec:certify_def}

\texttt{CERTIFY} is referred to throughout Chapters~\ref{ch:nofi} onward and is
defined here.

\begin{definition}[\texttt{CERTIFY} {\normalfont (\texttt{instr\_\allowbreak{}certify}, VMStep.v)}]
$\texttt{CERTIFY}(\Delta\mu)$ has the step semantics
\begin{equation}
\mathrm{pc} \mapsto \mathrm{pc}+1,\qquad
\mu \mapsto \mu + S(\Delta\mu),\qquad
\mathrm{cert} \mapsto \mathrm{true},
\end{equation}
with every other state component unchanged. Its cost is $S(\Delta\mu) \geq 1$,
which is the concrete instance of A2 (Section~\ref{sec:axiom_scheme}) for this
ISA. It is the \emph{only} rule in \texttt{VMStep.v} that sets
$\mathrm{cert} := \mathrm{true}$.
\end{definition}

\begin{remark}[\texttt{CERTIFY} is unguarded; what that does and does not mean]
\label{rem:certify_unguarded}
The rule fires for every state and every $\Delta\mu$: it carries no
precondition, takes no proposition as an argument, and checks nothing. Setting
the certification bit is therefore always \emph{available}, at a price of at
least $1$. Conversely, the instruction that does perform a decidable check on
data, namely \texttt{LASSERT}, which traps on failure, leaves both the
certification CSR and $\mathrm{cert}$ unchanged.

The honest description of the mechanism is therefore \emph{metering}, not
verification: the step rule guarantees that flipping the bit is priced, not
that the bit is only set when some claim holds. Statements of the form ``this
machine makes stamping \emph{verified} impossible'' overstate it; what is
impossible is stamping it \emph{for free}.

The ISA does contain genuinely guarded instructions.
\texttt{MORPH\_\allowbreak{}ASSERT} requires the named morphism to exist, and
\texttt{CHSH\_\allowbreak{}LASSERT} runs \texttt{column\_\allowbreak{}contractive\_\allowbreak{}check\_\allowbreak{}witness} on
actual witness data and traps on failure, with a soundness bridge to PSD. A
substrate in which the certification bit may be set only when an invariant
holds would let A2 be \emph{derived} (the cost of setting the bit becomes the
cost of establishing the invariant) rather than stipulated. That is an open
direction, not a present result; see Remark~\ref{rem:a2_stipulated}.
\end{remark}

\chapter{Operational Semantics}

\section{Transition Function}
The transition $\delta: \Sspace \times \Ispace \to \Sspace$ is defined by:
\begin{enumerate}
    \item $\mathit{PC}' = \mathit{PC} + 1$ for non-control-flow instructions; JUMP, JNEZ, CALL, and RET set PC to the target address; HALT stops execution.
    \item $\mu' = \mu + \Delta\mu_i$.
    \item $\mathit{err}' = \mathit{err} \lor \mathrm{fail}(S, i)$.
\end{enumerate}

\section{Traces and Execution}
\begin{definition}[Trace]
A trace $\tau = [i_0, \dots, i_k] \in \Ispace^*$ is a finite sequence of instructions.
\end{definition}

\begin{definition}[Execution]
$\mathrm{Run}([], S) = S$; \quad $\mathrm{Run}(i::\tau', S) = \mathrm{Run}(\tau', \delta(S,i))$.
\end{definition}

\begin{theorem}[Unique descended values, \textbf{(S)}]
Let $V(t)$ be evaluation of a VM instruction list from \texttt{init\_\allowbreak{}state}, and let $T(t)$ be target evaluation from a fixed basepoint. Then
\[
\bigl(\forall t,u,\ V(t)=V(u)\Rightarrow T(t)=T(u)\bigr)
\]
is equivalent to each reachable VM state having exactly one target value represented by its traces. This is \texttt{trace\_\allowbreak{}descent\_\allowbreak{}unique\_\allowbreak{}value\_\allowbreak{}iff}.
\end{theorem}

\begin{theorem}[Constructive reachable simulation, \textbf{(S)}]
Supply a function $r$ selecting an instruction list for each VM state, satisfying $V(r(V(t)))=V(t)$ for every trace $t$. A map preserving the basepoint, steps on reachable states, and certification exists exactly when target evaluation respects equal VM outcomes and target certification agrees with VM certification on every trace. The map is $s\mapsto T(r(s))$ and is pointwise unique on reachable states. This is \texttt{reachable\_\allowbreak{}simulation\_\allowbreak{}exists\_\allowbreak{}iff} together with \texttt{reachable\_\allowbreak{}simulation\_\allowbreak{}unique}. No selector construction or extension of the step law to unreachable states is asserted.
\end{theorem}

The certification quotient gives an explicit inhabited target. Conversely, the trace-history target satisfies A2 and certification agreement but fails descent: the empty trace and \texttt{JUMP 0 0} reach the same VM state while leaving distinct histories.

\section{Selected Instruction Semantics}

\textbf{PNEW}: Normalize region; if already present, idempotent; else add $(R_{\mathrm{norm}}, \emptyset)$ to $G$.

\textbf{REVEAL}: Update CSRs with checksum of certificate $\mathit{cert}$. Cost: $\mathit{bits} + S(\Delta\mu)$.

\textbf{LASSERT}: Read formula from \texttt{vm\_\allowbreak{}mem}[\texttt{R[formula\_\allowbreak{}addr\_\allowbreak{}reg]}] (encoded-unit count \texttt{formula\_\allowbreak{}len}, eight concrete bits per unit) and witness package from \texttt{vm\_\allowbreak{}mem}[\texttt{R[cert\_\allowbreak{}addr\_\allowbreak{}reg]}]. \texttt{cert\_\allowbreak{}kind = true} (SAT): verify using a satisfying model plus a falsifying countermodel (dual-witness). \texttt{cert\_\allowbreak{}kind = false} (UNSAT): the kernel step always fails at this path, \texttt{lassert\_\allowbreak{}check\_\allowbreak{}ok} returns \texttt{false} unconditionally; LRAT checking is not invoked in the step semantics (documented gap in \texttt{VMStep.v}). If valid, the constraint is certified; if invalid, set $\mathit{err} \gets \mathrm{true}$. Note: in the current hardware-aligned semantics, LASSERT validates the constraint but does not directly mutate the partition graph's axiom set, cert-channel activation happens through the subsequent CERTIFY or MORPH\_\allowbreak{}ASSERT. Cost: $\mathtt{formula\_\allowbreak{}len} \times 8 + S(\mathtt{mu\_\allowbreak{}delta})$.

%=============================================================================
% PART II: CORE THEOREMS
%=============================================================================

\chapter{No Free Insight}
\label{ch:nofi}

The central theorem: structural insight is never free.

\section{Definitions}

\begin{definition}[Receipt Predicate]
$P: \mathcal{O} \to \B$ is a computable function on observation histories.
\end{definition}

\begin{definition}[Strength Ordering]
$P_1 \le P_2$ iff $\forall o.\; P_1(o)=1 \implies P_2(o)=1$.
$P_1 < P_2$ iff $P_1 \le P_2$ and $\exists o.\; P_2(o)=1 \wedge P_1(o)=0$.
\end{definition}

\begin{definition}[Certification]
Trace $\tau$ \textbf{certifies} $P$ iff execution succeeds ($\mathit{err}=0$), the supra-certificate flag is set, and the output satisfies $P$.
\end{definition}

\begin{definition}[Structure Addition]
$\mathrm{HasStructureAddition}(\tau)$ is true iff $\tau$ contains a $\texttt{REVEAL}$, $\texttt{LASSERT}$, or equivalent operation.
\end{definition}

\section{The Theorem}

\begin{theorem}[No Free Insight {\normalfont (NoFreeInsight\_\allowbreak{}Theorem.v)}]
\label{thm:nfi}
Let $P_{\mathrm{strong}} < P_{\mathrm{weak}}$. If a trace $\tau$ starting from a clean state $s_0$ certifies $P_{\mathrm{strong}}$, then:
\[ \mathrm{Certified}(\tau, P_{\mathrm{strong}}) \implies \mathrm{HasStructureAddition}(\tau) \]
\end{theorem}

\begin{proof}[Proof Sketch]
The proof is parametric: \texttt{NoFreeInsight\_\allowbreak{}Theorem.v} proves the theorem for any system satisfying the \texttt{NO\_\allowbreak{}FREE\_\allowbreak{}INSIGHT\_\allowbreak{}SYSTEM} module type interface. The interface requires four contract obligations (not Coq \texttt{Axiom} declarations, they are proved separately for the kernel instantiation): (1)~$\mu$-monotonicity, (2)~a certification specification linking \texttt{certifies} to execution, (3)~a no-free-insight contract stating that strict strengthening requires structure addition, and (4)~that the system has a notion of clean start. The theorem file itself contains zero \texttt{Axiom}, zero \texttt{Admitted}. The contract obligations are discharged by the kernel instantiation.
\end{proof}

\begin{corollary}[Cost of Insight]
Gaining insight (strengthening the predicate) implies $\Delta\mu > 0$.
\end{corollary}

\section{$\mu$-Ledger Necessity}

\begin{theorem}[Observation adequacy, \textbf{(S)}]
For observation $o:S\to V$ and query $q:S\to A$, the existence of a decoder $d:V\to A$ satisfying $d(o(s))=q(s)$ implies
\[
o(s)=o(t)\ \Longrightarrow\ q(s)=q(t).
\]
Given a section $r:V\to S$ with $o(r(v))=v$, this condition is also sufficient: $d=q\circ r$. These are \texttt{decoding\_\allowbreak{}requires\_\allowbreak{}fiber\_\allowbreak{}constancy} and \texttt{selected\_\allowbreak{}representatives\_\allowbreak{}give\_\allowbreak{}decoder}.
\end{theorem}

\begin{theorem}[$\mu$-Ledger Necessity and Complete Independence Classification {\normalfont (NecessityOfMuLedger.v, kernel/NecessityAbstract.v)}, \textbf{(S)}]
\label{thm:mu_ledger_necessity}
Five independence results are formally proven, zero Admitted:
\begin{enumerate}
\item $\mu \perp (\texttt{mem}, \texttt{regs}, \texttt{pc})$: no strict-classical oracle recovers $\mu$ \hfill \texttt{turing\_\allowbreak{}ram\_\allowbreak{}mu\_\allowbreak{}necessity}
\item $\texttt{certified} \perp (\texttt{mem}, \texttt{regs}, \texttt{pc})$: no strict-classical oracle recovers certification \hfill \texttt{turing\_\allowbreak{}ram\_\allowbreak{}cert\_\allowbreak{}necessity}
\item $\texttt{certified} \perp (\texttt{mem}, \texttt{regs}, \texttt{pc}, \mu)$: knowing $\mu$ does not determine certification \hfill \texttt{cost\_\allowbreak{}model\_\allowbreak{}cert\_\allowbreak{}necessity}
\item $\mu \perp (\texttt{mem}, \texttt{regs}, \texttt{pc}, \texttt{certified})$: knowing certification does not determine $\mu$ \hfill \texttt{cert\_\allowbreak{}model\_\allowbreak{}mu\_\allowbreak{}necessity}
\item $\texttt{vm\_\allowbreak{}graph} \perp (\texttt{mem}, \texttt{regs}, \texttt{pc}, \mu, \texttt{certified})$: the categorical layer is independent of the full $\mu$-ledger shadow \hfill \texttt{graph\_\allowbreak{}not\_\allowbreak{}recoverable\_\allowbreak{}from\_\allowbreak{}P\_\allowbreak{}full}
\end{enumerate}
\textbf{Three-component independence} (\texttt{thiele\_\allowbreak{}state\_\allowbreak{}three\_\allowbreak{}component\_\allowbreak{}independence}): $\mu$, \texttt{vm\_\allowbreak{}certified}, and \texttt{vm\_\allowbreak{}graph} are the three provably independent non-classical state components. Each is irrecoverable from any combination of the others plus $(\texttt{mem}, \texttt{regs}, \texttt{pc})$.

\textbf{Minimality} (\texttt{P\_\allowbreak{}full\_\allowbreak{}complete\_\allowbreak{}neither\_\allowbreak{}mu\_\allowbreak{}nor\_\allowbreak{}cert\_\allowbreak{}droppable}): $(\texttt{mem}, \texttt{regs}, \texttt{pc}, \mu, \texttt{certified})$ is the minimal complete extension of classical state for $\mu$ and certification. Graph structure is a third independent axis.

\textbf{Universality} (\texttt{mu\_\allowbreak{}ledger\_\allowbreak{}necessity\_\allowbreak{}universal}): The separation holds from \emph{every} VMState, not only from crafted zero-state witnesses. From any reachable state, \texttt{CERTIFY 0} and \texttt{PNEW [] 0} produce identical strict classical shadows but different $\mu$. Certification also differs when the source state is uncertified.

\textbf{Trace closure} (\texttt{po1\_\allowbreak{}trace\_\allowbreak{}necessity}): For any common prefix program of any length, the two witness extensions maintain equal strict shadows but strictly different $\mu$. No accumulation of classical computation closes the gap.
\end{theorem}

\begin{proof}[Proof Sketch]
Each separation uses explicit witnesses where equal projected inputs force equal oracle outputs, while the witness states require unequal outputs, a contradiction. Strict shadow equality: both \texttt{CERTIFY 0} and \texttt{PNEW [] 0} advance \texttt{pc} by 1 and preserve \texttt{mem}/\texttt{regs}. $\mu$ divergence: CERTIFY charges $S(0)=1$; PNEW charges 0. Graph separation uses \texttt{categorical\_\allowbreak{}separation} from \texttt{PartitionSeparation.v}: two states with identical $P_{\text{full}}$ but different \texttt{pg\_\allowbreak{}morphisms}. Universality follows because \texttt{vm\_\allowbreak{}apply\_\allowbreak{}certify\_\allowbreak{}strict\_\allowbreak{}shadow} and \texttt{vm\_\allowbreak{}apply\_\allowbreak{}pnew\_\allowbreak{}strict\_\allowbreak{}shadow} hold for arbitrary source states.
\end{proof}

\begin{theorem}[Latent $\mu$: Trace Cost for a Fixed Schedule {\normalfont (NecessityOfMuLedger.v §7)}, \textbf{(S)}]
\label{thm:latent_mu}
The $\mu$-cost accumulated by any computation is a property of the instruction sequence, not of any tracking system. Four results, all zero Admitted:

\begin{enumerate}
\item \texttt{shadow\_\allowbreak{}mu\_\allowbreak{}is\_\allowbreak{}computation\_\allowbreak{}intrinsic}: For any program \texttt{prog} and any starting state $s$:
\[
(\texttt{run\_\allowbreak{}instrs}\ s\ \texttt{prog}).\mu = s.\mu + \texttt{trace\_\allowbreak{}total\_\allowbreak{}cost}(\texttt{prog})
\]
The increment is determined entirely by the instruction sequence, independent of vm\_\allowbreak{}graph, vm\_\allowbreak{}regs, vm\_\allowbreak{}mem, vm\_\allowbreak{}certified, and the starting $\mu$.

\item \texttt{shadow\_\allowbreak{}mu\_\allowbreak{}delta\_\allowbreak{}universal}: Two machines running the same program from any two starting states accumulate the same additional $\mu$. The cost is written in the program before execution begins.

\item \texttt{shadow\_\allowbreak{}mu\_\allowbreak{}unique\_\allowbreak{}accounting}: Any instruction-consistent accounting system $M$ satisfying $M(s) = 0$ at the starting state assigns exactly \texttt{trace\_\allowbreak{}total\_\allowbreak{}cost} to the result. Replace the Thiele Machine's $\mu$-counter with any other honest accounting and you get the same number.

\item \texttt{shadow\_\allowbreak{}mu\_\allowbreak{}inevitable}: Any program containing an instruction with positive declared cost accumulates strictly positive $\mu$ from any starting state. This conclusion requires a positive-cost instruction.
\end{enumerate}

\textbf{Interpretation.} Given a trace and this cost schedule, its sum is determined even if the observer does not retain the ledger. This does not establish a physical bill for every classical computation. A full encoding may retain the ledger; the irrecoverability theorem concerns the specified projection.
\end{theorem}

\section{Partition Refinement and No Free Insight}

\begin{theorem}[Partition Ops Are Free {\normalfont (kernel/PartitionRefinementNoFI.v)}, \textbf{(S)}]
PNEW, PSPLIT, and PMERGE are not cert-setters and can carry $\Delta\mu = 0$:
\begin{itemize}
    \item \texttt{partition\_\allowbreak{}structural\_\allowbreak{}ops\_\allowbreak{}not\_\allowbreak{}cert\_\allowbreak{}setters}: None of the three partition operations touch \texttt{csr\_\allowbreak{}cert\_\allowbreak{}addr} or \texttt{vm\_\allowbreak{}certified}.
    \item \texttt{partition\_\allowbreak{}structural\_\allowbreak{}ops\_\allowbreak{}can\_\allowbreak{}be\_\allowbreak{}free}: All three can execute with $\Delta\mu = 0$.
    \item \texttt{partition\_\allowbreak{}structural\_\allowbreak{}trace\_\allowbreak{}cannot\_\allowbreak{}certify}: A trace containing only PNEW, PSPLIT, and PMERGE (all at $\Delta\mu = 0$) cannot reach a certified state from an uncertified start.
\end{itemize}
Building partition structure costs nothing by itself. Certifying structural claims about that partition is what costs.
\end{theorem}

\begin{theorem}[Partition Free But Certification Non-Free {\normalfont (kernel/PartitionRefinementNoFI.v)}, \textbf{(S)}]
\label{thm:partition_nfi}
\texttt{partition\_\allowbreak{}free\_\allowbreak{}but\_\allowbreak{}certification\_\allowbreak{}nonfree}: There exist traces where all partition-structural instructions carry $\Delta\mu = 0$ (so $\mu_{\mathrm{final}} = \mu_{\mathrm{init}} = 0$) yet the trace ends with a partition graph that a subsequent MORPH\_\allowbreak{}ASSERT or CERTIFY can certify, but only at $\Delta\mu \ge 1$. Partition structure can be built at zero $\mu$-cost; certified partition knowledge cannot be acquired for free. Zero Admitted.
\end{theorem}

\begin{theorem}[Partition Refinement Non-Free {\normalfont (kernel/PartitionRefinementNoFI.v)}, \textbf{(S)}]
\texttt{partition\_\allowbreak{}refinement\_\allowbreak{}nonfree}: Any trace that starts with no cert evidence and ends with a certified claim about a partition structure paid at least 1 in $\mu$. This is the NoFI law specialized to partition refinement. Zero Admitted.
\end{theorem}

\chapter{The $\mu$-Chaitin Theorem}

A Chaitin-style incompleteness theorem denominated in $\mu$-currency.

\begin{theorem}[$\mu$-Chaitin {\normalfont (MuChaitinTheory\_\allowbreak{}Theorem.v)}]
For any formal theory system $T$ satisfying the $\mu$-Chaitin interface, the number of bits $k$ the theory can certify is bounded by:
\[ k \le |T| + c \]
where $|T|$ is the description length of the theory and $c$ is a fixed overhead constant. The proof chain: $k \le \mathrm{payload} \le \mu\text{-info} \le \mathrm{budget} \le |T| + c$.
\end{theorem}
\chapter{Computational Universality and Subsumption}
\section{Strict Containment}

\begin{theorem}[Turing $\subsetneq$ Thiele {\normalfont (kernel/Subsumption.v)}]
Classical Turing computation is \textbf{properly extended} by sighted Thiele computation (both are Turing-complete, but Thiele carries additional partition and $\mu$-accounting structure):
\[ \exists\, p.\; \mathrm{sighted}(p) \wedge \neg\mathrm{turing}(p) \]
There exist programs that are Thiele-computable (using partition structure and $\mu$-accounting) but are not classical Turing programs. The extra structure is genuine new content.
\end{theorem}
\section{Oracle Cost Lower Bound}
\begin{theorem}[Oracle pricing self-consistency (conceptual; not separately formalized)]
Inside the No Free Insight framework, the universal lower bound on certification cost (\texttt{universal\_\allowbreak{}nfi\_\allowbreak{}any\_\allowbreak{}substrate} in \texttt{coq/\allowbreak{}kernel/\allowbreak{}nfi/\allowbreak{}UniversalCertificationCost.v}) immediately implies that any sound oracle satisfying the one-step rule \texttt{cs\_\allowbreak{}cert\_\allowbreak{}costs} (cost $\ge 1$ on any uncertified-to-certified transition) charges $\ge n$ for $n$ certifying queries. There is no separate oracle-impossibility file in the codebase; the oracle case is a direct instantiation of the universal theorem at a \texttt{CertificationSystem} whose state encodes oracle responses. The physical justification for charging $\mu > 0$ on oracle queries comes from NoFI: gaining certified information about undecidable propositions has positive $\mu$-cost.
\end{theorem}
\section{$\mu$-budget predicates}

\begin{definition}[$\mu$-budget predicates {\normalfont (kernel/MuComplexity.v)}]
Let $p_{\mathrm{fuel}}, p_{\mu} : \N \to \N$ be polynomial bounds and $\mathrm{size} : \mathrm{VMState} \to \N$ a size function.
\begin{itemize}
    \item A decision problem $f$ is \textbf{\texttt{mu\_\allowbreak{}budget\_\allowbreak{}decidable}} (with respect to $p_{\mathrm{fuel}}$ and $p_{\mu}$) iff a solver trace exists for every input: a trace of length $\le p_{\mathrm{fuel}}(\mathrm{size}(s))$ using at most $p_{\mu}(\mathrm{size}(s))$ units of $\mu$-cost that runs without error.
    \item A decision problem $f$ is \textbf{\texttt{mu\_\allowbreak{}budget\_\allowbreak{}verifiable}} iff every positive instance has a certificate trace verifiable within the same fuel and $\mu$-cost bounds, ending with \texttt{vm\_\allowbreak{}certified = true}.
\end{itemize}
\end{definition}

\begin{theorem}[Zero-$\mu$ classical solvers are $\mu$-budget decidable {\normalfont (kernel/MuComplexity.v)}, \textbf{(S)}]
\texttt{classical\_\allowbreak{}mu\_\allowbreak{}budget\_\allowbreak{}decidable}: any decision problem solved by a zero-$\mu$ classical program within polynomial fuel is \texttt{mu\_\allowbreak{}budget\_\allowbreak{}decidable} with $\mu$-bound 0. Classical computation is a degenerate case of $\mu$-computation. Zero Admitted.
\end{theorem}

\begin{remark}
\texttt{mu\_\allowbreak{}budget\_\allowbreak{}decidable} and \texttt{mu\_\allowbreak{}budget\_\allowbreak{}verifiable} are predicates classifying decision problems by their polynomial-$\mu$ envelope, not classical complexity classes. The Coq file states decidable-implies-verifiable as an explicit premise; it proves the zero-$\mu$ solver inclusion above. It does not prove a class separation. The factored-search witness in \texttt{StructuralAdvantage.v} gives a concrete program-level $\mu$-cost advantage (paying $\mu = 18$ to certify a decomposition and reduce search from $O(N^2)$ to $O(2N)$), but that is not a class separation. The predicates are deliberately not named \texttt{in\_\allowbreak{}muP} and \texttt{in\_\allowbreak{}muNP}: those names evoke the classical P-versus-NP question, which this file does not address.
\end{remark}

\begin{theorem}[$\mu$-Cost Hierarchy {\normalfont (kernel/MuHierarchyTheorem.v)}, \textbf{(S)}]
\label{thm:mu-hierarchy}
For every $k \geq 1$, the $\mu$-dimension admits a strict hierarchy:
\begin{enumerate}
  \item \textbf{Achievability}: there exists a trace from \texttt{init\_\allowbreak{}state} with $\mu$-cost exactly $k$ that is \emph{level-$k$ certified}, the run executes a \texttt{CERTIFY}~$\delta$ instruction with $S(\delta) \geq k$ and ends with \texttt{vm\_\allowbreak{}certified = true}.
  \item \textbf{Lower bound}: any level-$k$ certified trace from \texttt{init\_\allowbreak{}state} must have spent at least $k$ units of $\mu$-cost.
\end{enumerate}
Formally (\texttt{mu\_\allowbreak{}hierarchy\_\allowbreak{}theorem}): for every $k \geq 1$, the witness trace $[\texttt{instr\_\allowbreak{}certify}(k{-}1)]$ (cost $= S(k{-}1) = k$) executes a level-$k$ \texttt{CERTIFY} from \texttt{init\_\allowbreak{}state}; conversely, the trace-side predicate \texttt{level\_\allowbreak{}k\_\allowbreak{}certified} forces $\texttt{trace\_\allowbreak{}mu\_\allowbreak{}cost} \geq k$ via \texttt{level\_\allowbreak{}k\_\allowbreak{}certification\_\allowbreak{}cost\_\allowbreak{}floor}, since the \texttt{CERTIFY}'s declared cost appears in the executed-instruction ledger and the ledger sum equals the trace cost (\texttt{run\_\allowbreak{}vm\_\allowbreak{}mu\_\allowbreak{}conservation}).

\textbf{Corollary} (\texttt{mu\_\allowbreak{}hierarchy\_\allowbreak{}no\_\allowbreak{}upper\_\allowbreak{}bound}): no fixed $\mu$-budget suffices for all levels. The $\mu$-dimension is unbounded and non-collapsible. Zero Admitted.
\end{theorem}

\begin{remark}[$\mu$-Hierarchy as a P-vs-NP analogue, not a P-vs-NP result]
Theorem~\ref{thm:mu-hierarchy} is the $\mu$-analogue of the classical time hierarchy theorem. It is not a statement about classical P versus NP, which this project does not touch. In the certified-structure setting it makes one thing precise: \emph{finding} a certificate of structural complexity $k$ costs $\mu \geq k$, while \emph{verifying} it (given the trace) costs $0$ extra $\mu$. The hierarchy is strict and unbounded, with zero \texttt{Admitted}, no oracle, and no diagonalization. The lower-bound proof is trace-side: a level-$k$ certified run executed a specific \texttt{CERTIFY}~$\delta$ with $S(\delta) \geq k$, and \texttt{executed\_\allowbreak{}instruction\_\allowbreak{}cost\_\allowbreak{}recorded} plus \texttt{ledger\_\allowbreak{}sum\_\allowbreak{}contains\_\allowbreak{}lower\_\allowbreak{}bound} place that cost in the ledger sum, which \texttt{run\_\allowbreak{}vm\_\allowbreak{}mu\_\allowbreak{}conservation} equates with the trace's $\mu$-cost.
\end{remark}

\section{Intrinsic level hierarchy}

\begin{definition}[Intrinsic level {\normalfont (coq/IntrinsicLevelHierarchy.v)}]
A final state $s$ is \emph{intrinsically certifiable at level $\geq k$}, written $\texttt{level\_\allowbreak{}intrinsic\_\allowbreak{}at\_\allowbreak{}least}(k, s)$, iff every $\texttt{run\_\allowbreak{}vm}$ trace from \texttt{init\_\allowbreak{}state} that ends at $s$ with \texttt{vm\_\allowbreak{}certified = true} executes at least $k$ cert-setter instructions:
\begin{multline*}
\forall \mathit{fuel}\,\mathit{trace}.\;
\texttt{run\_\allowbreak{}vm}(\mathit{fuel},\mathit{trace},\texttt{init\_\allowbreak{}state})=s \\
{}\land\; s.\texttt{vm\_\allowbreak{}certified}=\texttt{true}
\;\Rightarrow\; \texttt{cert\_\allowbreak{}setter\_\allowbreak{}executions}(\mathit{fuel},\mathit{trace},\texttt{init\_\allowbreak{}state})\geq k.
\end{multline*}
The level is a property of the \emph{final state}, not of any particular trace reaching it.
\end{definition}

\begin{theorem}[Intrinsic level forces $\mu$ {\normalfont (coq/IntrinsicLevelHierarchy.v::level\_\allowbreak{}intrinsic\_\allowbreak{}forces\_\allowbreak{}mu)}]
If $\texttt{level\_\allowbreak{}intrinsic\_\allowbreak{}at\_\allowbreak{}least}(k, s)$ holds and $s.\texttt{vm\_\allowbreak{}certified} = \texttt{true}$, then for every trace reaching $s$ from \texttt{init\_\allowbreak{}state}, $s.\texttt{vm\_\allowbreak{}mu} \geq k$. Zero Admitted.
\end{theorem}

\begin{theorem}[Strict separation {\normalfont (coq/IntrinsicLevelHierarchy.v::level\_\allowbreak{}k\_\allowbreak{}plus\_\allowbreak{}1\_\allowbreak{}forces\_\allowbreak{}more\_\allowbreak{}than\_\allowbreak{}k\_\allowbreak{}cert\_\allowbreak{}events)}]
\label{thm:intrinsic-strict-separation}
For every $k \geq 0$, if $\texttt{level\_\allowbreak{}intrinsic\_\allowbreak{}at\_\allowbreak{}least}(k+1, s)$ holds and $s.\texttt{vm\_\allowbreak{}certified} = \texttt{true}$, then every trace reaching $s$ executes strictly more than $k$ cert-setter instructions. Zero Admitted.
\end{theorem}

\begin{remark}[Relation to Theorem~\ref{thm:mu-hierarchy}]
Theorem~\ref{thm:intrinsic-strict-separation} is the state-level counterpart of the trace-level $\mu$-hierarchy. The trace-level version (\texttt{mu\_\allowbreak{}hierarchy\_\allowbreak{}theorem}, \texttt{level\_\allowbreak{}k\_\allowbreak{}certification\_\allowbreak{}cost\_\allowbreak{}floor}) defines level on the trace, a trace is level-$k$ certified iff it executes a \texttt{CERTIFY}~$\delta$ with $S(\delta) \geq k$. The state-level version defines level on the final state, every trace reaching it requires at least $k$ cert-setter executions. The two are complementary entry points to the same hierarchy: the trace-level theorem gives the achievability witness ($[\texttt{instr\_\allowbreak{}certify}(k{-}1)]$ from \texttt{init\_\allowbreak{}state}); the state-level theorem gives the lower bound on $\mu$ for any state intrinsically requiring $k$ cert-events. Both are present-tense in the kernel; both close under the global context.
\end{remark}

\section{Verifier model and the bare-setting impossibility}

\begin{definition}[Bare verifier {\normalfont (coq/VerifierModel.v)}]
The bare-setting transcript type is $\texttt{BareTranscript} := \texttt{list StrictClassicalState}$. A bare verification problem is a record \texttt{BareVerificationProblem} carrying
\begin{itemize}
    \item \texttt{bare\_\allowbreak{}honest\_\allowbreak{}claim} : $\texttt{VMState} \to \texttt{Prop}$;
    \item \texttt{bare\_\allowbreak{}explains} : $\texttt{VMState} \to \texttt{BareTranscript} \to \texttt{Prop}$ (the transcript-honestly-witnessed-by-state relation);
    \item \texttt{bare\_\allowbreak{}cheap\_\allowbreak{}budget}, \texttt{bare\_\allowbreak{}recompute\_\allowbreak{}cost} : $\N$.
\end{itemize}
A bare verifier is a record \texttt{BareVerifier} with $\texttt{bv\_\allowbreak{}decide} : \texttt{BareTranscript} \to \texttt{bool}$ and $\texttt{bv\_\allowbreak{}cost} : \texttt{BareTranscript} \to \N$.
\end{definition}

\begin{definition}[Soundness, completeness, cheapness {\normalfont (coq/VerifierModel.v)}]
For a problem $P$ and verifier $V$:
\begin{itemize}
    \item \texttt{bare\_\allowbreak{}sound}$(P, V) := \forall t.\; V(t) = \texttt{true} \Rightarrow \forall s.\; \texttt{bare\_\allowbreak{}explains}(P, s, t) \Rightarrow \texttt{bare\_\allowbreak{}honest\_\allowbreak{}claim}(P, s)$.
    \item \texttt{bare\_\allowbreak{}complete}$(P, V) := \forall s\, t.\; \texttt{bare\_\allowbreak{}honest\_\allowbreak{}claim}(P, s) \Rightarrow \texttt{bare\_\allowbreak{}explains}(P, s, t) \Rightarrow V(t) = \texttt{true}$.
    \item \texttt{bare\_\allowbreak{}cheap}$(P, V) := \forall t.\; \texttt{bv\_\allowbreak{}cost}(V, t) \leq \texttt{bare\_\allowbreak{}cheap\_\allowbreak{}budget}(P)$.
\end{itemize}
The soundness condition is \emph{strong}: the claim must hold for every state explaining the transcript, not merely for some.
\end{definition}

\begin{theorem}[Bare-setting impossibility {\normalfont (coq/VerifierImpossibility.v::bare\_\allowbreak{}setting\_\allowbreak{}no\_\allowbreak{}sound\_\allowbreak{}complete\_\allowbreak{}verifier)}]
\label{thm:bare-verifier-impossibility}
There is no bare verifier simultaneously \texttt{bare\_\allowbreak{}sound} and \texttt{bare\_\allowbreak{}complete} on the verification problem $\texttt{mu\_\allowbreak{}eq\_\allowbreak{}one\_\allowbreak{}problem}$ whose claim is $\texttt{vm\_\allowbreak{}mu}(s) = 1$ and whose \texttt{bare\_\allowbreak{}explains} relation holds for the two $\texttt{NecessityOfMuLedger}$ witness states against the shared strict-shadow trace.
\end{theorem}

\begin{proof}
$\texttt{po1\_\allowbreak{}state\_\allowbreak{}A}$ (the cost-1 \texttt{CERTIFY 0} result, $\mu = 1$) and $\texttt{po1\_\allowbreak{}state\_\allowbreak{}B}$ (the cost-0 \texttt{PNEW [] 0} result, $\mu = 0$) project to identical strict-shadow traces. Completeness applied to $A$'s honest run forces $V$ to accept the shared transcript. Strong soundness applied to the same accepted transcript at the explanation by $B$ forces $\texttt{vm\_\allowbreak{}mu}(B) = 1$, contradicting $\texttt{vm\_\allowbreak{}mu}(B) = 0$. \texttt{Print Assumptions} returns \texttt{Closed under the global context}.
\end{proof}

\begin{theorem}[Three sufficient escapes {\normalfont (coq/VerifierEscape\_\allowbreak{}\{Substrate,Hardness,Interaction\}.v)}]
\label{thm:three_escapes}
Each kernel theorem below exhibits a constant-cost verifier that simultaneously satisfies soundness and completeness on the $\mu$-sensitive claim by augmenting the transcript with non-classical structure:
\begin{itemize}
    \item \texttt{substrate\_\allowbreak{}escape\_\allowbreak{}succeeds}, transcript is $\texttt{list VMState}$; verifier reads $\mu$ directly.
    \item \texttt{hardness\_\allowbreak{}escape\_\allowbreak{}succeeds}, transcript is $\texttt{BareTranscript} \times \texttt{bool}$ (commitment bit); verifier accepts the commitment under a parameterised \texttt{HardnessHypothesis} record.
    \item \texttt{interactive\_\allowbreak{}escape\_\allowbreak{}succeeds}, transcript is $\texttt{BareTranscript} \times \N$ (response); verifier checks the response against the claim.
\end{itemize}
Each closes under the global context with zero \texttt{Admitted}.
\end{theorem}

\begin{theorem}[Factorisation impossibility {\normalfont (coq/VerifierExhaustiveness.v::V\_\allowbreak{}does\_\allowbreak{}not\_\allowbreak{}factor\_\allowbreak{}through\_\allowbreak{}classical)}]
\label{thm:verifier-factorisation}
Let $T$ be any transcript type with a projection $\pi : T \to \texttt{BareTranscript}$. For any verifier $V : T \to \texttt{bool}$ simultaneously sound and complete on the $\mu$-sensitive claim (with the lifted-explains relation), $V$ does not factor through $\pi$: there exist $t_1, t_2 : T$ with $\pi(t_1) = \pi(t_2)$ and $V(t_1) \neq V(t_2)$.
\end{theorem}

\begin{remark}[The verifier corollary]
Sound and complete verification of $\mu$-sensitive claims requires the transcript to carry non-classical structural information, exactly the structure the projection $\texttt{strict\_\allowbreak{}shadow}$ drops. The three escapes above are concrete realisations; Theorem~\ref{thm:verifier-factorisation} says no verifier on the bare classical projection alone succeeds, regardless of how rich its transcript is otherwise. The substrate channel is the option the structural axis makes available; hardness and interaction are the two routes classical cryptography and complexity already used.
\end{remark}

%=============================================================================
% PART IV: SYMMETRY AND CONSERVATION
%=============================================================================

\chapter{$\mu$-Conservation and Receipt Integrity}

\section{The $\mu$-Ledger Identity}

\begin{theorem}[$\mu$-Ledger Conservation {\normalfont (MuLedgerConservation.v)}]
For any trace $\tau = [i_0, \dots, i_k]$:
\[ \mu_{\mathrm{final}} = \mu_{\mathrm{init}} + \sum_{j=0}^{k} \mathrm{cost}(i_j) \]
Every instruction increases $\mu$ by exactly its declared cost. No axioms, no admits.
\end{theorem}

\begin{corollary}[Monotonicity]
$\mu$ never decreases: $\mu(s) \le \mu(\mathrm{Run}(\tau, s))$ for all $\tau$.
\end{corollary}

\begin{corollary}[Irreversibility Bound]
$\mathrm{irreversible\_\allowbreak{}bits}(\tau) \le \mu_{\mathrm{final}} - \mu_{\mathrm{init}}$.
\end{corollary}

\begin{theorem}[$\mu$ Decomposition {\normalfont (MuLedgerConservation.v)}]
$\mu_{\mathrm{total}} = \mu_{\mathrm{blind}} + \mu_{\mathrm{sighted}}$ for any partition of the cost into blind (reversible) and sighted (structural) components.
\end{theorem}

\section{Receipt Integrity}

\begin{definition}[Receipt]
A receipt $r = (\mathrm{step}, \mathrm{instr}, \mu_{\mathrm{pre}}, \mu_{\mathrm{post}}, h_{\mathrm{pre}}, h_{\mathrm{post}})$ records one transition.
\end{definition}

\begin{theorem}[Receipt Integrity {\normalfont (ReceiptIntegrity.v)}]
A valid receipt or execution attestation proves $\mu_{\mathrm{final}} = \mu_{\mathrm{init}} + \sum \mathrm{costs}$ for the attested run. Any receipt claiming $\Delta\mu \neq \mathrm{cost}(\mathrm{instr})$ fails validation.
\end{theorem}

\begin{theorem}[Non-Forgeability {\normalfont (ReceiptIntegrity.v)}]
Forged receipts (claiming incorrect $\Delta\mu$) fail validation. Overflow values ($\mu > 2^{31}-1$) are rejected at the receipt layer.
\end{theorem}

\chapter{Gauge Symmetry and Noether's Theorem}

\section{The Partial $\Z$-Action (Gauge Group)}

\begin{definition}[Gauge Shift]
For $\delta \in \Z$, the gauge shift $\sigma_\delta: \Sspace \to \Sspace$ maps $S \mapsto S[\mu \gets \mathrm{toNat}(\mu + \delta)]$, preserving all other fields.

Note the truncation. Since $\mu$ ranges over $\N$, the Coq definition
\texttt{z\_\allowbreak{}gauge\_\allowbreak{}shift} computes \texttt{Z.to\_\allowbreak{}nat (Z.of\_\allowbreak{}nat vm\_\allowbreak{}mu + delta)},
and \texttt{Z.to\_\allowbreak{}nat} \emph{clamps negatives to zero}. So $\sigma_\delta$ is a
total function on states but is \emph{not} injective, and the shifts do not
form a $\Z$-action on all of $\Sspace$.
\end{definition}

\begin{theorem}[Partial Group Action {\normalfont (KernelNoether.v)}]
The gauge shifts form a \emph{partial} $\Z$-action on states; equivalently, a
genuine $\Z$-action on the sub-poset of states where no intermediate shift
underflows:
\begin{enumerate}
    \item \textbf{Identity}: $\sigma_0(S) = S$. Unconditional
      (\texttt{z\_\allowbreak{}action\_\allowbreak{}identity}).
    \item \textbf{Composition}: $\sigma_a(\sigma_b(S)) = \sigma_{a+b}(S)$
      \emph{provided} $\mu(S) + b \ge 0$ \emph{and} $\mu(S) + b + a \ge 0$
      (\texttt{z\_\allowbreak{}action\_\allowbreak{}composition}, which carries both hypotheses).
    \item \textbf{Inverse}: $\sigma_{-n}(\sigma_n(S)) = S$ provided
      $\mu(S) + n \ge 0$ (\texttt{z\_\allowbreak{}action\_\allowbreak{}inverse}).
\end{enumerate}
\end{theorem}

\begin{remark}[Why composition needs its guard]
The guard on clause~2 is not decoration. Take $\mu(S) = 0$, $b = -1$, $a = +1$:
$\sigma_b$ clamps to $\mu = 0$, so $\sigma_a(\sigma_b(S))$ has $\mu = 1$;
whereas $\sigma_{a+b}(S) = \sigma_0(S) = S$ has $\mu = 0$. The two sides differ.
Composition fails wherever an intermediate shift underflows, which is why the
mechanised theorem carries \emph{two} nonnegativity hypotheses: one for the
intermediate state, one for the final state. Drop either and the counterexample
above goes through.
\end{remark}

\section{Gauge Invariance}

\begin{definition}[Observable]
$\mathrm{Obs}(S)$ extracts the partition region structure, ignoring $\mu$.
\end{definition}

\begin{theorem}[Gauge Invariance {\normalfont (KernelNoether.v)}]
$\mathrm{Obs}(\sigma_\delta(S)) = \mathrm{Obs}(S)$. Shifting $\mu$ by a constant does not affect observables.
\end{theorem}

\begin{theorem}[Noether's Theorem {\normalfont (\texttt{coq/\allowbreak{}kernel/\allowbreak{}curvature/\allowbreak{}KernelNoether.v})}]
\leavevmode
\begin{enumerate}
    \item \textbf{Forward}: States identical except for $\mu$ lie on the same gauge orbit. (This is direct from the definition: if all other fields match, the $\mu$ difference is the gauge shift.)
    \item \textbf{Gauge Invariance}: \texttt{z\_\allowbreak{}gauge\_\allowbreak{}invariance} proves that shifting $\mu$ does not change \texttt{Observable\_\allowbreak{}partition}. This is definitional: \texttt{z\_\allowbreak{}gauge\_\allowbreak{}shift} only modifies the \texttt{vm\_\allowbreak{}mu} field, and \texttt{Observable\_\allowbreak{}partition} reads only \texttt{vm\_\allowbreak{}graph}.
    \item \textbf{Trace Preservation}: Gauge-equivalent states produce identical observable traces (same labels, same $\mu$-costs) for any finite horizon. Proven as \texttt{vm\_\allowbreak{}step\_\allowbreak{}orbit\_\allowbreak{}equiv}.
\end{enumerate}
The $\mu$-ledger is a gauge degree of freedom. Its \emph{absolute} value is unobservable; only \emph{differences} $\Delta\mu$ are physical. The ``Noether'' label is a deliberate analogy: the result follows from the record structure (separate fields for $\mu$ and graph), not from a deep symmetry principle.

\begin{remark}[Gauge freedom and absolute $\mu$-claims]
\label{rem:gauge_vs_absolute}
Gauge freedom has to be squared with the results that assert absolute values, most directly the verifier argument, which needs $\mathrm{vm\_\allowbreak{}mu}(s) = 1$ to be a
well-posed claim, and the $\mu$-hierarchy, which needs ``cost exactly $k$''. If
only differences carried meaning, neither statement would have content.

Every absolute $\mu$-claim in this document is \emph{relative to}
\texttt{init\_\allowbreak{}state}, where $\mu = 0$ by definition. Read
$\mathrm{vm\_\allowbreak{}mu}(s) = 1$ as $\Delta\mu = 1$ along a run from
\texttt{init\_\allowbreak{}state} to $s$; $\mu$-conservation
(\texttt{run\_\allowbreak{}vm\_\allowbreak{}mu\_\allowbreak{}conservation}) is what licenses the reading, since it
makes the final ledger equal to the trace cost.

The scope is worth naming precisely: $\sigma_\delta$ is defined on
\emph{arbitrary} states, not on orbits rooted at \texttt{init\_\allowbreak{}state}, so the
gauge group as formalised is larger than the equivalence the absolute claims
need. Nothing in the development quotients by it. The honest statement is that
the theorems live on \emph{reachable} states, where the base point pins the
gauge, and the gauge chapter describes a symmetry of the state \emph{type}
rather than of the reachable dynamics. Note also that
\texttt{vm\_\allowbreak{}step\_\allowbreak{}mu\_\allowbreak{}monotonic} makes $\mu$ non-decreasing along runs, which
already breaks the $\Z$-symmetry on trajectories: the dynamics single out a
direction the gauge action does not respect.
\end{remark}
\end{theorem}

\begin{theorem}[$\mu$-Monotonicity {\normalfont (KernelNoether.v)}]
$\mathrm{vm\_\allowbreak{}step}(S, i, S') \implies \mu(S) \le \mu(S')$. The $\mu$-ledger never decreases under the dynamics. This is true by construction: instruction costs are natural numbers ($\ge 0$), so adding them to the ledger cannot decrease it. The proof examines each step constructor and confirms the arithmetic.
\end{theorem}

%=============================================================================
% PART V: CONDITIONAL QUANTUM ANALOGS
%=============================================================================

\begin{remark}[Part V Scope: Conditional Results and Named Axioms]
The chapters in this part present results conditional on named physical axioms (linearity in the Bloch parameter for the Born rule, the second law for Landauer, and the \texttt{A\_\allowbreak{}QM} bridge for quantum realizability). They are marked \textbf{(C)} or \textbf{(R)} throughout. The primary unconditional quantum results, the classical CHSH bound, the algebraic Tsirelson bound, and the biconditional \texttt{column\_\allowbreak{}contractive\_\allowbreak{}iff\_\allowbreak{}quantum\_\allowbreak{}realizable}, appear earlier: the classical CHSH bound in Chapter~\ref{ch:bell}, the algebraic Tsirelson bound in Chapter~\ref{ch:tsirelson}, and the biconditional in Chapter~\ref{ch:biconditional}. Those are pure algebra, no physics axioms.

The results in this part are not primary claims of the project. They are explorations of what follows \emph{if} you add specific physical axioms. They show how far the machine reaches without claiming any physics the proofs do not support.
\end{remark}

\chapter{The Born Rule}

\begin{definition}[Bloch Sphere Probabilities]
For a state with purity vector $(x,y,z)$:
\[ P(|0\rangle) = \frac{1+z}{2}, \qquad P(|1\rangle) = \frac{1-z}{2} \]
\end{definition}

\begin{definition}[Measurement $\mu$-Cost]
\[ \mathrm{Cost}(x,y,z) = \frac{1 - (x^2 + y^2 + z^2)}{2} \]
This is the linear entropy. For pure states ($r^2=1$), cost $= 0$. For mixed states ($r^2 < 1$), cost $> 0$.
\end{definition}

\begin{theorem}[Uniqueness of the Born Rule {\normalfont (BornRule.v)}]
The Born rule $P(|0\rangle) = (1+z)/2$ is the \textbf{unique} probability assignment satisfying:
\begin{enumerate}
    \item Non-negativity and normalization: $P \ge 0$, $P(|0\rangle) + P(|1\rangle) = 1$.
    \item Eigenstate consistency: $P(0,0,1) = 1$, $P(0,0,-1) = 0$.
    \item Linearity in $z$: $P(x,y,z) = az + b$ for some $a,b$.
\end{enumerate}
Proof: Boundary conditions force $a = 1/2$, $b = 1/2$. No other solution exists.

\textbf{Note:} The Coq theorem also lists a \texttt{receipt\_\allowbreak{}mu\_\allowbreak{}consistent} hypothesis, but it is \textbf{unused} in the proof. The hypothesis $\texttt{receipt\_\allowbreak{}mu\_\allowbreak{}consistent}~P$ unfolds to ``measurement cost $\ge 0$ for valid Bloch vectors,'' which is directly from the definition true for all states and does not constrain $P$. The linearity assumption (item~3) does all the work, it assumes the functional form $P = az+b$ and solves a 2-equation system. This is less ``deriving the Born rule from $\mu$-accounting'' and more ``the Born rule is the unique linear rule consistent with eigenstates.''
\end{theorem}

\chapter{Unitarity}

\begin{definition}[Evolution]
An evolution $E$ maps Bloch vectors $(x,y,z) \mapsto (x',y',z')$ with associated $\mu$-cost.
\end{definition}

\begin{theorem}[Unitarity from Zero Cost {\normalfont (Unitarity.v)}]
\leavevmode
\begin{enumerate}
    \item $\mu = 0 \implies r'^2 \ge r^2$ (purity cannot decrease at zero cost). Proven as \texttt{zero\_\allowbreak{}cost\_\allowbreak{}preserves\_\allowbreak{}purity}.
    \item Non-unitary evolution ($r'^2 < r^2$ for some state) requires $\mu > 0$. Proven as \texttt{nonunitary\_\allowbreak{}requires\_\allowbreak{}mu}: the proof specializes the conservation hypothesis at the witness, then \texttt{lra}.
    \item $\mu = 0 \Rightarrow$ no info loss (one direction proven). The converse and the full equivalence $\mu = 0 \Leftrightarrow$ unitary $\Leftrightarrow$ reversible are \emph{stated} in the file but not proven as theorems: they appear as \texttt{Definition} declarations (propositional constants) rather than proven \texttt{Theorem}s.
\end{enumerate}
\end{theorem}

\begin{theorem}[Lindblad Requires $\mu$ {\normalfont (Unitarity.v)}]
Lindblad-type dissipation at rate $\gamma > 0$ requires $\mu \ge \gamma$. Decoherence is paid for.
\end{theorem}

\begin{theorem}[CPTP Structure {\normalfont (Unitarity.v)}]
Positivity + trace preservation $\implies$ the evolution is CPTP (completely positive, trace-preserving).
\end{theorem}

\chapter{Purification}

\begin{theorem}[Purification Principle {\normalfont (Purification.v)}]
Every mixed state $(x,y,z)$ with $r^2 = x^2+y^2+z^2 \le 1$ admits eigenvalues $\lambda_1, \lambda_2$ with:
\begin{itemize}
    \item $\lambda_1 + \lambda_2 = 1$, \quad $0 \le \lambda_i \le 1$.
    \item $(\lambda_1 - \lambda_2)^2 = r^2$.
    \item Construction: $\lambda_1 = (1 + r)/2$, $\lambda_2 = (1 - r)/2$.
\end{itemize}
Pure states ($r=1$) need no external reference (deficit $= 0$). The maximally mixed state ($r=0$) has deficit $= 1$.
\end{theorem}

\chapter{No-Cloning}

\begin{theorem}[No-Cloning from $\mu$-Conservation {\normalfont (NoCloning.v)}]
Let $I = x^2 + y^2 + z^2$ be the information content (purity). For a cloning operation with outputs of information $I_1, I_2$ and $\mu$-cost $\mu$:
\[ I_1 + I_2 \le I + \mu \]
\end{theorem}

\begin{theorem}[Approximate Cloning {\normalfont (NoCloning.v)}]
For fidelities $f_1, f_2$: $f_1 + f_2 \le 1 + \mu/I$. At zero cost: $f_1 + f_2 \le 1$.
\end{theorem}

\begin{theorem}[No Deletion {\normalfont (NoCloning.v)}]
Perfect deletion also requires $\mu \ge I$ (dual of no-cloning).
\end{theorem}

\chapter{The Tsirelson Bound}
\label{ch:tsirelson}

\begin{remark}[Epistemological Status: \textbf{(S)} algebraic, no physics axioms]
The active-corpus derivation lives in \texttt{coq/\allowbreak{}kernel/\allowbreak{}quantum/\allowbreak{}TsirelsonGeneral.v} (\texttt{tsirelson\_\allowbreak{}from\_\allowbreak{}row\_\allowbreak{}bounds}) and \texttt{coq/\allowbreak{}kernel/\allowbreak{}category/\allowbreak{}AlgebraicCoherence.v} (\texttt{tsirelson\_\allowbreak{}bound\_\allowbreak{}tight}); both prove the bound algebraically without encoding $2\sqrt{2}$ into any definition. \texttt{TsirelsonGeneral.v} derives $S^2 \le 8$ from pure algebra: a sum-of-squares identity proves each CHSH row contributes $\le 2(e_1^2 + e_2^2)$, then Cauchy--Schwarz bounds the sum by $4\sum e_i^2 \le 8$. \texttt{TsirelsonFromAlgebra.v} provides a self-contained algebraic derivation with achievability witness ($e = \pm 1/\sqrt{2}$, giving $S = 2\sqrt{2}$) and a verified rational bound ($\sqrt{8} < 5657/2000$). Zero admits in both files.
\end{remark}

\begin{theorem}[Tsirelson Upper Bound for the Zero-Marginal NPA Model {\normalfont (TsirelsonQuantumModel.v, TsirelsonFromAlgebra.v)}, \textbf{(S)}]
If a trace is quantum-realizable in the zero-marginal NPA model (i.e., is slice-coherent / column-contractive), then $|S_{\mathrm{CHSH}}| \le 2\sqrt{2}$.

The algebraic derivation (\texttt{TsirelsonFromAlgebra.v}) proves $S^2 \le 8$ for all algebraically coherent correlators satisfying the three column-contractivity conditions, using a sum-of-squares identity and psatz. No physics premises. Pure polynomial arithmetic.
\end{theorem}

\chapter{Witness Insight Taxonomy}

\begin{theorem}[CHSH Trial Is Tier 0 {\normalfont (kernel/WitnessInsightGeneral.v)}, \textbf{(S)}]
\texttt{chsh\_\allowbreak{}trial\_\allowbreak{}is\_\allowbreak{}tier0}: The CHSH\_\allowbreak{}TRIAL instruction never touches either certification channel:
\begin{itemize}
    \item \texttt{chsh\_\allowbreak{}trial\_\allowbreak{}not\_\allowbreak{}cert\_\allowbreak{}addr\_\allowbreak{}setter}: CHSH\_\allowbreak{}TRIAL does not change \texttt{csr\_\allowbreak{}cert\_\allowbreak{}addr}.
    \item \texttt{chsh\_\allowbreak{}trial\_\allowbreak{}preserves\_\allowbreak{}vm\_\allowbreak{}certified}: CHSH\_\allowbreak{}TRIAL does not change \texttt{vm\_\allowbreak{}certified}.
\end{itemize}
Recording CHSH trial outcomes is provably free of structural cost. The counters go up; the cert channels are untouched. Cost: 0. Zero Admitted.
\end{theorem}

\begin{theorem}[Three-Tier Observation Taxonomy {\normalfont (kernel/WitnessInsightGeneral.v)}, \textbf{(S)}]
\label{thm:taxonomy}
\texttt{witness\_\allowbreak{}insight\_\allowbreak{}complete\_\allowbreak{}taxonomy}: Every observation the machine can make falls into exactly one of three tiers:

\begin{description}
    \item[\textbf{Tier 0}] Raw observation (always free): Records data without activating any certification channel. CHSH\_\allowbreak{}TRIAL is the canonical example. Cost: $\Delta\mu = 0$ is possible; no cert channel is touched.
    \item[\textbf{Tier 1}] Certified structural insight (always costs $\ge 1$): Any instruction that activates \texttt{csr\_\allowbreak{}cert\_\allowbreak{}addr} or \texttt{vm\_\allowbreak{}certified} is in this tier. Proven: \texttt{certified\_\allowbreak{}witness\_\allowbreak{}insight\_\allowbreak{}nonfree} shows $\mu$ increases by at least 1.
    \item[\textbf{Tier 2}] Certified non-local witness insight (always costs $\ge 1$): A certified execution whose WitnessCount statistics show $S > 2$. \texttt{nonlocal\_\allowbreak{}witness\_\allowbreak{}insight\_\allowbreak{}nonfree}: getting certified with non-local evidence requires $\mu \ge 1$.
\end{description}

\texttt{witness\_\allowbreak{}insight\_\allowbreak{}nonfree\_\allowbreak{}general}: For any instruction in Tier 1 or Tier 2, $\mu$ increases by at least 1. The only free observations are Tier 0. Zero Admitted.
\end{theorem}

\chapter{Bell's Theorem in Morphism Coupling Language}
\label{ch:bell}

\begin{definition}[Setting-Outcome Coupling {\normalfont (kernel/CHSHCouplingBridge.v)}]
The \texttt{chsh\_\allowbreak{}setting\_\allowbreak{}outcome\_\allowbreak{}coupling} maps a WitnessCount record to a coupling (list of $(nat \times nat)$ pairs). Each pair records a (setting-pair-index, outcome-type) observation. A coupling is \textbf{separable} (functional) when it maps each setting to at most one outcome type; \textbf{non-separable} otherwise.
\end{definition}

\begin{theorem}[Locally Consistent Strategy Gives Separable Coupling {\normalfont (kernel/CHSHCouplingBridge.v)}, \textbf{(S)}]
\texttt{locally\_\allowbreak{}consistent\_\allowbreak{}gives\_\allowbreak{}separable\_\allowbreak{}coupling}: Any WitnessCount configuration produced by a locally deterministic strategy $(a_0, a_1, b_0, b_1) \in \{0,1\}^4$ gives a separable (functional) coupling. Zero Admitted.
\end{theorem}

\begin{theorem}[Local Coupling Bound {\normalfont (kernel/CHSHCouplingBridge.v)}, \textbf{(S)}]
\texttt{locally\_\allowbreak{}consistent\_\allowbreak{}classical\_\allowbreak{}bound}: Any locally factorizable morphism coupling satisfying \texttt{WCLocallyConsistent} has $|S| \le 2$. Bell's theorem in the morphism coupling language: the separation between local and non-local correlations is not just a statistical fact but a structural property of the coupling predicate. Zero Admitted.
\end{theorem}

\begin{theorem}[CHSH Violation Rules Out Local Couplings {\normalfont (kernel/CHSHCouplingBridge.v)}, \textbf{(S)}]
\texttt{chsh\_\allowbreak{}violation\_\allowbreak{}rules\_\allowbreak{}out\_\allowbreak{}locally\_\allowbreak{}factorizable\_\allowbreak{}coupling}: If WitnessCount statistics satisfy $S > 2$, no locally factorizable morphism coupling is consistent with those counts. Contrapositively: non-local correlations cannot be produced by any local morphism coupling. The no-go is stated in the coupling language, not just the statistical layer. Zero Admitted.
\end{theorem}

\begin{remark}
These results connect the categorical morphism layer directly to Bell nonlocality. A claim that two modules are entangled (via MORPH\_\allowbreak{}ASSERT on a non-separable coupling) cannot be locally faked: the NoFI law charges at least 1 for the MORPH\_\allowbreak{}ASSERT, and the coupling structure itself rules out any local hidden strategy that would produce $S > 2$ for free.
\end{remark}

\chapter{Slice Coherence and NPA Realizability: The Biconditional}
\label{ch:biconditional}

This algebraic result has no project-local axioms or admitted lemmas; its assumption receipt includes standard classical-real and functional-extensionality axioms.

\begin{definition}[Column Contractivity (Slice Coherence)]
\section{The CHSH statistic $S$ and its witness counts}
\label{sec:chsh_statistic}

The central quantum object of this part is defined here, since it is used
throughout without ever having been given as a function of machine state.

\begin{definition}[Correlator from witness counts {\normalfont (CHSHStatisticalBridge.v)}]
A VM state carries a \texttt{WitnessCounts} record in $\N^8$, holding for each
setting pair $(x,y) \in \{0,1\}^2$ the tallies
$\mathrm{same}_{xy}$ and $\mathrm{diff}_{xy}$. Writing
$N_{xy} = \mathrm{same}_{xy} + \mathrm{diff}_{xy}$, the empirical correlator is
\begin{equation}
E_{xy} =
\begin{cases}
\dfrac{\mathrm{same}_{xy} - \mathrm{diff}_{xy}}{N_{xy}} & N_{xy} > 0,\\[2ex]
0 & N_{xy} = 0 .
\end{cases}
\end{equation}
The $N_{xy}=0$ branch is a definitional convention (no trials, no evidence), not
a claim that uncorrelated is the right default.
\end{definition}

\begin{definition}[CHSH statistic {\normalfont (\texttt{chsh\_\allowbreak{}S}, TsirelsonFromMu.v)}]
\begin{equation}
S = E_{00} + E_{01} + E_{10} - E_{11} \in \mathbb{R} .
\end{equation}
Note the sign on the last term, and note that in
\texttt{CHSHStatisticalBridge.v} the same statistic is assembled with the
$(1,1)$ arguments reversed, giving $E^{*}_{11} = (\mathrm{diff}_{11} -
\mathrm{same}_{11})/N_{11}$, so that all four terms may be added. The two
presentations agree.

The classical bound is $|S| \leq 2$ and the Tsirelson bound is
$|S| \leq 2\sqrt{2}$. Both are proved here as \emph{algebraic} facts with no
physics axioms. Where the development states the Tsirelson bound in squared
form it appears as $S^2 \leq 8$, which is the same bound: $\sqrt{8} = 2\sqrt2$.
\end{definition}

A correlation tuple $(E_{00}, E_{01}, E_{10}, E_{11})$ is \textbf{slice-coherent} (equivalently, column-contractive; the predicate tests slice membership rather than truthfulness, and the name records that) if:
\begin{align}
1 - E_{00}^2 - E_{10}^2 &\geq 0 \quad \text{(condition 1)} \\
1 - E_{01}^2 - E_{11}^2 &\geq 0 \quad \text{(condition 2)} \\
(1-E_{00}^2-E_{10}^2)(1-E_{01}^2-E_{11}^2) &\geq (E_{00}E_{01}+E_{10}E_{11})^2 \quad \text{(condition 3)}
\end{align}
These are exactly the PSD conditions on the $5 \times 5$ zero-marginal NPA moment matrix.
\end{definition}

\begin{theorem}[Slice Coherence $\iff$ NPA Realizability {\normalfont (\texttt{coq/\allowbreak{}kernel/\allowbreak{}quantum/\allowbreak{}QuantumPartitionPSD.v}, corollary \texttt{column\_\allowbreak{}contractive\_\allowbreak{}iff\_\allowbreak{}quantum\_\allowbreak{}realizable})}, \textbf{(S)}]
\label{thm:thiele_honest_iff_qr}
\texttt{column\_\allowbreak{}contractive\_\allowbreak{}iff\_\allowbreak{}quantum\_\allowbreak{}realizable}:
\[
\text{slice\_\allowbreak{}coherent}(\mathbf{E}) \iff \text{quantum\_\allowbreak{}realizable}(\mathbf{E})
\]
where \emph{slice\_\allowbreak{}coherent} means the three column-contractivity conditions hold, and \emph{quantum\_\allowbreak{}realizable} means the $5\times5$ zero-marginal NPA moment matrix is positive semidefinite.

\textbf{Forward direction} (\texttt{zero\_\allowbreak{}marginal\_\allowbreak{}npa\_\allowbreak{}column\_\allowbreak{}contractive\_\allowbreak{}implies\_\allowbreak{}psd} in \texttt{coq/\allowbreak{}kernel/\allowbreak{}nfi/\allowbreak{}MuLedgerQuantumBridge.v}): column-contractive $\Rightarrow$ NPA PSD. Proved by constructing the Cholesky decomposition of the $5\times5$ matrix from the three conditions.

\textbf{Reverse direction} (\texttt{npa\_\allowbreak{}psd\_\allowbreak{}implies\_\allowbreak{}column\_\allowbreak{}contractive} in \texttt{coq/\allowbreak{}kernel/\allowbreak{}quantum/\allowbreak{}QuantumPartitionPSD.v}): NPA PSD $\Rightarrow$ column-contractive. Proved by specializing the PSD condition at the vector $(0, -(E_{00}y_0+E_{01}y_1), -(E_{10}y_0+E_{11}y_1), y_0, y_1)$, obtaining a non-negative quadratic in $y_0, y_1$, then closing via \texttt{quadratic\_\allowbreak{}nonneg\_\allowbreak{}discriminant}.

This biconditional has no project-local axioms; its assumption report includes standard-library real-number and functional-extensionality assumptions. No physics assumptions. Pure polynomial arithmetic. The zero-marginal NPA framework is the orthogonal-observables slice of the quantum elliptope ($\rho_{AA}=\rho_{BB}=0$), not the full elliptope.

Key consequences:
\begin{itemize}
\item PR box $(E_{00}=E_{01}=E_{10}=1, E_{11}=-1)$ fails condition 1 immediately ($1-1-1=-1<0$). Direct arithmetic; no separate theorem needed.
\item \texttt{state\_\allowbreak{}column\_\allowbreak{}contractive\_\allowbreak{}implies\_\allowbreak{}tsirelson} (\texttt{coq/\allowbreak{}kernel/\allowbreak{}nfi/\allowbreak{}MuLedgerQuantumBridge.v}) directly proves slice-coherent $\Rightarrow$ $|S|^2 \leq 8$ (equivalently $|S| \leq 2\sqrt{2}$); its proof composes \texttt{state\_\allowbreak{}column\_\allowbreak{}contractive\_\allowbreak{}implies\_\allowbreak{}quantum\_\allowbreak{}gram} with \texttt{quantum\_\allowbreak{}realizable\_\allowbreak{}implies\_\allowbreak{}tsirelson\_\allowbreak{}bound}. The Cauchy--Schwarz inner step (PSD $\Rightarrow$ row bounds $\Rightarrow$ $|S|^2 \leq 8$) is \texttt{tsirelson\_\allowbreak{}from\_\allowbreak{}row\_\allowbreak{}bounds} in \texttt{coq/\allowbreak{}kernel/\allowbreak{}quantum/\allowbreak{}TsirelsonGeneral.v}. The contrapositive ($|S| > 2\sqrt{2}$ $\Rightarrow$ not slice-coherent) follows immediately. The converse is false and is not claimed.
\item Turing-classical correlations $(1,0,1,0)$ satisfy $|S|=2$ but fail condition 1: $1-1-1=-1<0$. A Turing machine can output this correlation; the slice gate rejects it as slice-incoherent (the full-elliptope gate of the next section accepts it, the point is inside the elliptope, outside the zero-marginal slice). Direct arithmetic from \texttt{column\_\allowbreak{}contractive\_\allowbreak{}iff\_\allowbreak{}quantum\_\allowbreak{}realizable}.
\end{itemize}
Zero Admitted. Zero named hypotheses.
\end{theorem}

\section{The Elliptope Completion and Its Gate}

The biconditional above characterizes the zero-marginal slice. The full correlator quantum set is characterized in \texttt{coq/\allowbreak{}kernel/\allowbreak{}quantum/\allowbreak{}ElliptopeCompletion.v} as an existential completion: $(E_{00}, E_{01}, E_{10}, E_{11})$ is \textbf{elliptope-realizable} (\texttt{elliptope\_\allowbreak{}realizable}) iff some assignment of the cross moments $x = \langle A_0 A_1\rangle$, $y = \langle B_0 B_1\rangle$ makes the completed $5\times5$ moment matrix PSD.

\begin{theorem}[Elliptope completion, kernel-tier]
All of the following hold with \texttt{Print Assumptions} reporting only the stdlib classical-reals and functional-extensionality families:
\begin{itemize}
\item \texttt{deterministic\_\allowbreak{}strategy\_\allowbreak{}elliptope}: every deterministic sign strategy is elliptope-realizable, witnessed by its own cross moments (rank-one Gram argument).
\item \texttt{elliptope\_\allowbreak{}finite\_\allowbreak{}mixture} / \texttt{lhv\_\allowbreak{}mixture\_\allowbreak{}elliptope}: the set is closed under arbitrary finite sub-convex mixtures; every LHV correlator is inside.
\item \texttt{elliptope\_\allowbreak{}tsirelson}: every member satisfies $S^2 \leq 8$, via Cauchy--Schwarz inside the PSD form (\texttt{psd\_\allowbreak{}cauchy\_\allowbreak{}schwarz}); row-norm bounds provably fail on the full set.
\item \texttt{pr\_\allowbreak{}box\_\allowbreak{}not\_\allowbreak{}elliptope}: the PR box admits no completion.
\item \texttt{beyond\_\allowbreak{}classical\_\allowbreak{}elliptope}: $S = 12/5 > 2$ is attained inside, so the classical inclusion is strict.
\end{itemize}
\end{theorem}

The decidable gate (\texttt{coq/\allowbreak{}kernel/\allowbreak{}quantum/\allowbreak{}ElliptopeGate.v}) is the $\mathbb{Z}$-arithmetic check an opcode can branch on: a fraction-free Sylvester branch over supplied completion buckets (strict interior) and a rational $LDL^{\mathsf T}$-certificate branch reaching singular and boundary completions, with soundness (\texttt{elliptope\_\allowbreak{}check\_\allowbreak{}full\_\allowbreak{}sound}) into \texttt{elliptope\_\allowbreak{}realizable}, computed acceptances of the $\mu = 0$ tightness witness, $(1,0,1,0)$, and the on-Tsirelson-curve point $(3/5,4/5,4/5,-3/5)$ (\texttt{gate\_\allowbreak{}accepts\_\allowbreak{}pythagorean\_\allowbreak{}boundary}), and the universal refusal \texttt{elliptope\_\allowbreak{}full\_\allowbreak{}gate\_\allowbreak{}never\_\allowbreak{}accepts\_\allowbreak{}pr\_\allowbreak{}box}. Scope: correlator level; the only points beyond any exact integer check are those whose every PSD completion is irrational; the runtime opcode binding is future engineering.

\section{The Pointer-Observable Criterion}

\texttt{coq/\allowbreak{}kernel/\allowbreak{}frontier/\allowbreak{}PointerObservable.v} formalizes the schema behind the metering conjecture. An \textbf{ecosystem} (\texttt{Ecosystem}) is a state type carrying $n$ observers, each holding a fragment map over states; an event $E$ \textbf{redundantly proliferates} (\texttt{redundantly\_\allowbreak{}proliferating}) iff every observer's fragment records $E$; and $E$ is the \textbf{unique pointer among} a designated list of rival events (\texttt{unique\_\allowbreak{}pointer\_\allowbreak{}among}) iff $E$ proliferates and no listed rival does. Uniqueness is relative to the named rivals on purpose, absolute uniqueness is false for free. A toy replicated-ledger instance shows the definitions are satisfiable and discriminating: the certification event is the unique pointer among the toy's two meterable events (\texttt{toy\_\allowbreak{}cert\_\allowbreak{}unique\_\allowbreak{}pointer}), a sanity instance only.

\texttt{PointerObservableReductions.v} instantiates the schema at the five deployed metering disciplines of the reductions tier, proof-of-stake finality, gas metering, TEE attestation, certificate transparency, and proof-carrying verification. Each discipline's metered event is proved the unique pointer among its named rivals (\texttt{PoS\_\allowbreak{}unique\_\allowbreak{}pointer}, \texttt{Gas\_\allowbreak{}unique\_\allowbreak{}pointer}, \texttt{TEE\_\allowbreak{}unique\_\allowbreak{}pointer}, \texttt{CT\_\allowbreak{}unique\_\allowbreak{}pointer}, \texttt{PCC\_\allowbreak{}unique\_\allowbreak{}pointer}), and the conjunction is exported as \texttt{five\_\allowbreak{}disciplines\_\allowbreak{}are\_\allowbreak{}pointers}, which closes under the global Coq context, no axioms at all, stdlib or otherwise.

Scope, stated plainly: these theorems verify that the five formalized ecosystems have the claimed pointer structure. They are evidence for the pointer-observable conjecture, that deployed metering disciplines meter exactly the event whose record proliferates, not a proof of it. The forced-event question stays open, and the file header says so.

\subsection{The adversarial search}
\label{sec:pointer_counterexamples}

Five confirmations settle nothing on their own. A criterion loose enough to accept anything accepts five things just as easily, and what makes a universality class convincing is never the confirmations. It is the reported absence of exceptions after somebody went looking. \texttt{coq/\allowbreak{}kernel/\allowbreak{}frontier/\allowbreak{}PointerObservableCounterexamples.v} is the looking, and it is aimed at the corner where the criterion is most likely to break rather than at systems resembling the five.

The claim put at risk is the strong reading:

\begin{quote}
\textbf{(PO-STRONG)} Every system that must resist forgery of a commitment has a commitment event that redundantly proliferates.
\end{quote}

A counterexample must satisfy three conditions fixed in advance: (C1) it must resist forgery of a commitment, in the sense that its designers treat an unearned commitment as the primary failure mode; (C2) it must be deployed and independently evolved, owing nothing to this development; and (C3) its commitment event must fail \texttt{redundantly\_\allowbreak{}proliferating}.

\textbf{(PO-STRONG) is false.} Three candidates refute it, each machine-checked:

\begin{itemize}
    \item \textbf{Deniable (designated-verifier) authentication} (\texttt{deniable\_\allowbreak{}authentication\_\allowbreak{}refutes\_\allowbreak{}strong\_\allowbreak{}criterion}). The scheme is engineered so a third party cannot be convinced even given the whole transcript, because the designated verifier could have produced it. Non-transferability is the product requirement.
    \item \textbf{Symmetric-key authentication} (\texttt{mac\_\allowbreak{}refutes\_\allowbreak{}strong\_\allowbreak{}criterion}). Existential unforgeability is the defining security notion of a MAC, and verification requires the shared secret, so no record reaches anyone else.
    \item \textbf{Object capabilities} (\texttt{capability\_\allowbreak{}refutes\_\allowbreak{}strong\_\allowbreak{}criterion}). Unforgeability here is unconditional, enforced by memory safety rather than by a cost or a hardness assumption. Nothing is paid, nothing is logged, and no observer outside the holder can tell.
\end{itemize}

A public-log control (\texttt{public\_\allowbreak{}log\_\allowbreak{}confirms}) checks that the search can return \emph{confirms} and is not an instrument that only reports refutations; its rival effort event still fails to proliferate, so the control discriminates rather than accepting everything.

The boundary these trace is sharp, and it is not the one (PO-STRONG) guessed at. A MAC and a digital signature both resist forgery; only the signature proliferates. Deniable authentication and certificate transparency both resist forgery; deniability actively prevents third-party conviction while transparency actively requires it. Object capabilities and proof-carrying code both make unearned authority impossible; the capability confines it to the holder while the proof travels to every consumer. In each pair the proliferating member is the one whose commitment must convince parties who were not present. The criterion that survives is therefore:

\begin{quote}
\textbf{(PO-PUBLIC)} A system whose commitment must convince third parties who did not witness its creation will proliferate records of the commitment event.
\end{quote}

The signature observer model makes the commitment visible in each observer fragment and proves proliferation for that model. Its correspondence to deployed signatures remains a modeling judgment. Public availability of a verifier does not imply that every observer retains a record.

The verifier constructions give three sufficient evidence interfaces under their respective premises. The hardness construction has the weaker soundness guarantee stated in its theorem. Neither these examples nor the factorization obstruction proves that only one route can involve metering.

Three statements travel together in loose prose and come apart under these candidates:

\begin{description}
    \item[(M1)] \textbf{Forgery resistance forces metering.} \textbf{False.} Signatures take the hardness route; object capabilities achieve unforgeability unconditionally through memory safety, with no cost, no record, and no assumption.
    \item[(M2)] \textbf{Forgery resistance forces record proliferation.} \textbf{False.} Deniable authentication, symmetric MACs, and capabilities all resist forgery while proliferating nothing, and in the first case that is the design goal.
    \item[(M3)] \textbf{Public commitments force record proliferation.} \textbf{Proposed.} The chosen observer models satisfy it. Actual record proliferation needs independent evidence; public verifiability alone is insufficient.
\end{description}

What remains of the metering claim, stated so a stranger can attack it: when a claim is $\mu$-sensitive, meaning it is not decidable from the classical projection, and the verifier can neither recheck it directly nor substitute a hardness assumption, the substrate route is what is left, and that route prices the commitment event. That is \texttt{V\_\allowbreak{}does\_\allowbreak{}not\_\allowbreak{}factor\_\allowbreak{}through\_\allowbreak{}classical} in different clothes. It is narrower than ``forgery resistance,'' and it is the version this development supports.

%=============================================================================
% PART VI: EMERGENT STRUCTURE
%=============================================================================

\chapter{Emergent Spacetime}

\section{The $\mu$-Metric}

\begin{definition}[$\mu$-Distance {\normalfont (MuGeometry.v)}]
The Coq formalization defines \texttt{mu\_\allowbreak{}distance} as the absolute difference of $\mu$-totals:
\[ d_\mu(S_1, S_2) = |\mu_{\mathrm{total}}(S_2) - \mu_{\mathrm{total}}(S_1)| \]
This is a one-dimensional projection, not a minimum over paths.
\end{definition}

\begin{theorem}[Pseudometric Properties {\normalfont (MuGeometry.v)}]
$d_\mu$ satisfies non-negativity, reflexivity ($d(S,S)=0$), symmetry, and triangle inequality, all inherited from absolute value on $\mathbb{Z}$. Note: $d_\mu$ is a \textbf{pseudometric}, not a metric: $d(S_1,S_2)=0$ does \emph{not} imply $S_1=S_2$ (many distinct states share the same $\mu$-total). The ``geometry'' is a 1D projection onto the $\mu$ counter.
\end{theorem}

\section{Causal Cones}

\begin{definition}[$\mu$-Reachability Cone {\normalfont (SpacetimeEmergence.v)}]
$C^+(S) = \{ S' \mid \exists \tau.\; S \xrightarrow{\tau} S' \wedge \mathrm{cost}(\tau) \le \mathrm{Budget} \}$.
\end{definition}

This is a structural analogy to causal cones: the set of reachable states given a finite $\mu$-budget plays a role analogous to a light cone. The analogy should not be over-read: there is no Lorentz invariance, no metric signature, and no boost transformations in the Coq formalization.

\section{No-Signaling}

\begin{theorem}[Observational Locality {\normalfont (SpacetimeEmergence.v)}]
If module $M_B$ is not in the target set of any instruction in a trace, then $M_B$'s observable region is unchanged after execution (\texttt{exec\_\allowbreak{}trace\_\allowbreak{}no\_\allowbreak{}signaling\_\allowbreak{}outside\_\allowbreak{}cone}).
\end{theorem}

\begin{remark}
This is a \textbf{data isolation} property: operations targeting module $A$ in the partition graph do not affect the lookup of module $B$. The proof is structural (list operations on disjoint keys), non-trivial mainly for \texttt{instr\_\allowbreak{}psplit}/\texttt{instr\_\allowbreak{}pmerge} which restructure the graph. The label ``no-signaling'' is an analogy. There is no spatial metric or speed limit between modules.
\end{remark}

\section{Discrete Geometry: The Gauss-Bonnet Result}
\label{sec:gauss_bonnet}

The partition graph is a discrete geometric object. One unconditional result falls out, requiring no physics axioms.

\begin{theorem}[Discrete Angle-Defect Identity (Planar Triangulations) {\normalfont (\texttt{coq/\allowbreak{}kernel/\allowbreak{}curvature/\allowbreak{}DiscreteGaussBonnet.v})}, \textbf{(S)}]
For any \texttt{well\_\allowbreak{}formed\_\allowbreak{}triangulated} partition graph $g$ in the equilateral angle model (each triangle corner contributes $\pi/3$):
\[
\mathrm{angle\_\allowbreak{}defect}(g) = 5\pi \cdot \chi(g)
\]
where $\chi(g) = V - E + F$ is the Euler characteristic (vertices minus edges plus faces), and $\mathrm{angle\_\allowbreak{}defect}(g)$ is the sum over all vertices of $2\pi$ minus the total incident angle. This is a fully proven combinatorics result about the discrete triangulation model. It is not the continuum Gauss-Bonnet theorem and it is not a derivation of angles from $\mu$-cost. Zero Admitted.
\end{theorem}

\begin{remark}[How this differs from classical Gauss-Bonnet]
\textbf{The factor $5\pi$ is not a typo and not the classical Gauss-Bonnet constant.} For a closed triangulated 2-manifold (like the surface of a tetrahedron), the standard derivation gives $D = 2\pi\chi$: every interior edge is shared by exactly two faces, so $3F = 2E$, and substituting that into $D = 2\pi V - \pi F$ together with $\chi = V - E + F$ yields $D = 2\pi\chi$. The Coq predicate \texttt{well\_\allowbreak{}formed\_\allowbreak{}triangulated} does \emph{not} include $3F = 2E$. It instead requires \texttt{satisfies\_\allowbreak{}boundary\_\allowbreak{}euler\_\allowbreak{}relation} ($B = 3\chi$, where $B$ is the boundary edge count) together with $E = I + B$ and $3F = 2I + B$. From these the file derives the combinatorial identity $3V = 5E - 6F$, which in turn forces $D = 5\pi\chi$.

That identity holds for \textbf{planar triangulations with boundary} (the canonical example: a planar disk with $V=7$, $E=15$, $F=9$, $\chi=1$, satisfying $3V = 21 = 5E - 6F$ and $D = 5\pi$). It \textbf{fails} for closed surfaces (the tetrahedron has $V=4$, $E=6$, $F=4$, but $3V = 12 \neq 6 = 5E - 6F$, so it is excluded from the predicate). The theorem is therefore not a discrete Gauss-Bonnet theorem in the classical sense, it is a different combinatorial identity for a different family of complexes. The two coexist: classical Gauss-Bonnet at $2\pi\chi$ for closed 2-manifolds, the file's identity at $5\pi\chi$ for planar disks satisfying its well-formedness predicate.

This is the full extent of the topology-curvature identity in the project.
\end{remark}

\subsection{Cross-link in MuGravity Vocabulary}

\begin{theorem}[Graph-Laplacian Sum-Zero Identity {\normalfont (\texttt{coq/\allowbreak{}kernel/\allowbreak{}frontier/\allowbreak{}F3\_\allowbreak{}MuLaplacianSum.v})}, \textbf{(S)}]
\label{thm:total_mu_laplacian_zero}
For every \texttt{VMState} $s$:
\[
\sum_{m \in \mathrm{pg\_\allowbreak{}modules}(\mathrm{vm\_\allowbreak{}graph}\,s)} \mathrm{mu\_\allowbreak{}laplacian}(s, m) = 0.
\]
Proof by edge-pair antisymmetry. \texttt{nat\_\allowbreak{}list\_\allowbreak{}disjoint} is symmetric, hence \texttt{modules\_\allowbreak{}adjacent\_\allowbreak{}by\_\allowbreak{}region} is symmetric, hence \texttt{edge\_\allowbreak{}weight} is symmetric; \texttt{mu\_\allowbreak{}gradient} is antisymmetric by definition; the double sum cancels under the Fubini swap. No well-formedness, no calibration. Standalone graph-theoretic fact.
\end{theorem}

\begin{theorem}[F3 Narrow Cross-link {\normalfont (\texttt{coq/\allowbreak{}kernel/\allowbreak{}frontier/\allowbreak{}F3\_\allowbreak{}PartitionTopologyCrossLink.v})}, \textbf{(S)}]
\label{thm:f3_narrow}
For every \texttt{VMState} $s$ at universal local calibration ($\forall m, \texttt{calibration\_\allowbreak{}residual}(s, m) = 0$):
\[
\sum_{m \in \mathrm{pg\_\allowbreak{}modules}(\mathrm{vm\_\allowbreak{}graph}\,s)} \mathrm{geometric\_\allowbreak{}angle\_\allowbreak{}defect}(s, m) = 0.
\]
The proof composes two independent kernel facts: \texttt{calibration\_\allowbreak{}residual\_\allowbreak{}zero\_\allowbreak{}iff} (geometry $\leftrightarrow$ $\mu$-ledger at every module) and \texttt{total\_\allowbreak{}mu\_\allowbreak{}laplacian\_\allowbreak{}zero} (Theorem~\ref{thm:total_mu_laplacian_zero}). Both ingredients are load-bearing, drop counter-instances \texttt{F3\_\allowbreak{}drop\_\allowbreak{}calibration\_\allowbreak{}breaks\_\allowbreak{}prediction} and \texttt{F3\_\allowbreak{}drop\_\allowbreak{}sum\_\allowbreak{}zero\_\allowbreak{}breaks\_\allowbreak{}prediction} retained.
\end{theorem}

\begin{remark}[Why this is not the topology version]
\label{rem:f3_no_topology}
The discrete angle-defect identity does not imply $\sum_m \mathrm{mu\_\allowbreak{}laplacian} = 5 \cdot \chi$. That inference would require identifying two distinct angle-defect quantities: \texttt{DiscreteGaussBonnet.angle\_\allowbreak{}defect} (degree $\times$ $\pi/3$, equilateral) and \texttt{MuGravity.geometric\_\allowbreak{}angle\_\allowbreak{}defect} (distance-weighted, $\pi \cdot d_{bc}/(1 + d_{ab}+d_{ac}+d_{bc})$). They are not equal in general: \texttt{strong\_\allowbreak{}bridge\_\allowbreak{}counterexample} (in the same file) formalises a degree-6 vertex at uniform inter-module distance $d \geq 1$ where the Gauss-Bonnet defect is $0$ and the MuGravity defect is $2\pi/(1 + 3d) \neq 0$. The cross-link makes no claim about the topology invariant. Discrete Gauss-Bonnet remains a kernel theorem; it is simply not load-bearing for this cross-link.

\textbf{Non-vacuity is open.} Symbolic search across small actually-triangulated graphs ($K_4$, $K_4 - e$, octahedron, prism, $K_5$, $B_3$, $K_6$, $K_7$) up to $N = 7$ returns no integer-mass calibrated solution; structural analysis identifies three failure modes (insufficient triangles, excess triangles, asymptotic-only target reachability). \texttt{module\_\allowbreak{}triangles} was tightened in \texttt{MuGravity.v} to require pairwise-adjacent neighbors (an actual graph triangle), eliminating tree-shaped pseudo-witnesses.
\end{remark}

\section{Four-Dimensional Discrete Riemannian Geometry (Reference Only, Not a Main Claim)}

\begin{remark}[This section is not about spacetime]
\label{rem:not_spacetime}
The metric defined below is diagonal with entries
$2\cdot|\mathrm{module\_\allowbreak{}structural\_\allowbreak{}mass}|$,
which are non-negative, so the signature is $(+,+,+,+)$: the form is
\emph{Riemannian}, not Lorentzian. There is no timelike direction, no light
cone, no causal structure, and no proper time. Calling a positive-definite
4-dimensional metric ``spacetime'' misdescribes the object even under a
reference-only label.

Read this section as: discrete Riemannian geometry on a 4-dimensional
simplicial complex built from partition-graph data. Every general-relativistic
word that appears below (Christoffel, Riemann, Einstein, stress-energy) is
being used for the corresponding \emph{Riemannian} construction. The vacuum
result remains a $0 = 0$ consistency check, as the scope warning says.
\end{remark}

\begin{remark}[Scope Warning]
The following subsections define discrete differential geometry objects (simplicial complex, metric tensor, Christoffel symbols, Riemann tensor, Einstein tensor) over partition graphs and prove results about the vacuum case $G_{\mu\nu} = 0 = T_{\mu\nu}$. This is a $0 = 0$ result: both sides are zero in the vacuum because mass is zero everywhere. It is \textbf{not} a derivation of Einstein's field equations from computation. The gravitational constant $G := 1/(8\pi)$ is a unit choice (a definition), not a derived quantity. No physical claim is being made. The geometry section's main result is the Discrete Gauss-Bonnet theorem (Section~\ref{sec:gauss_bonnet}). Everything in this section is marked reference-only and does not appear in the primary claim ledger.
\end{remark}

The previous sections treated spacetime as a 1D projection via the $\mu$-metric. The next subsections extend to full 4D geometry as a reference exploration. None of the results here are part of the core formal claims of the project.

\subsection{4D Simplicial Complex}

\begin{definition}[4D Simplicial Complex {\normalfont (FourDSimplicialComplex.v)}]
A \textbf{4D simplicial complex} $\mathcal{K}$ is a discrete approximation to 4D spacetime consisting of:
\begin{itemize}
    \item \textbf{Vertices} $V$: Computational modules (ModuleID)
    \item \textbf{Edges} $E \subseteq V \times V$: Connections between modules
    \item \textbf{Simplices}: 2-simplices (triangles), 3-simplices (tetrahedra), 4-simplices (4D cells)
\end{itemize}
The complex satisfies:
\begin{itemize}
    \item \textbf{Consistency}: Every face of a simplex is also in $\mathcal{K}$
    \item \textbf{Well-formedness}: Dimension 4 with locally Euclidean structure
\end{itemize}
\end{definition}

\subsection{Metric Tensor from $\mu$-Costs}

\begin{definition}[Computational Metric {\normalfont (MetricFromMuCosts.v, RiemannTensor4D.v)}]
The metric tensor at vertex $v$ with respect to vertices $w_1, w_2$ is:
\begin{equation}
g_{\mu\nu}(v; w_1, w_2) = \begin{cases}
\ell(v, w_1, w_2) & \text{if } \mu = \nu \\
0 & \text{if } \mu \neq \nu
\end{cases}
\end{equation}
where the edge length function is:
\begin{equation}
\ell(s, v_1, v_2) = 2 \cdot \text{mass}(v_1) \quad \text{if } v_1 = v_2
\end{equation}
and mass is derived from structural complexity: $\text{mass}(v) = |\text{module\_\allowbreak{}structural\_\allowbreak{}mass}(s, v)|$.
\end{definition}

\begin{remark}[The metric is specified only on the diagonal]
\label{rem:metric_underspecified}
$\ell$ is given only for $v_1 = v_2$. No value is stated for $v_1 \neq v_2$, so
as written $g$ is a partial specification. This matters downstream rather than
cosmetically: the Christoffel symbols require $\partial_\mu g_{\nu\sigma}$, and
the discrete derivative \emph{differences $g$ across neighbouring vertices}. It evaluates $g$ exactly where the definition above is silent. The
Coq development supplies a total function, so the mechanised objects are
well-defined; this document simply does not state which extension it uses.
Treat the equations below as schematic until the off-diagonal case is written
down here. Since the whole section is reference-only and the vacuum result is a
$0=0$ check, nothing in the claim ledger depends on the resolution.

Also note $g$ is positive semi-definite, so it degenerates wherever
$\mathrm{mass}(v) = 0$, including the entire vacuum configuration, which is
the case the theorems below actually evaluate. A degenerate metric has no
inverse, so ``raising an index'' is not available there.
\end{remark}

\begin{remark}
This metric is \textbf{position-dependent} for non-uniform mass distributions:
\begin{itemize}
    \item \textbf{Uniform mass} $m$: $g_{\mu\mu}(v; w, w) = 2m$ (constant) $\implies$ flat spacetime
    \item \textbf{Non-uniform mass}: $g_{\mu\mu}(v; w, w)$ varies with $w$ $\implies$ curved spacetime
\end{itemize}
In this model, mass distributions create position-dependent geometry: uniform mass gives flat spacetime, non-uniform mass gives curved spacetime.
\end{remark}

\subsection{Discrete Differential Geometry}

\begin{definition}[Discrete Derivative {\normalfont (RiemannTensor4D.v)}]
For a scalar field $f: V \to \mathbb{R}$ and direction $\mu$:
\begin{equation}
\partial_\mu f(v) = \frac{1}{|N_\mu(v)|} \sum_{w \in N_\mu(v)} \bigl(f(w) - f(v)\bigr)
\end{equation}
where $N_\mu(v)$ is the set of neighbors of $v$ in the $\mu$-direction.
\end{definition}

\begin{definition}[Christoffel Symbols {\normalfont (RiemannTensor4D.v)}]
The connection coefficients (Christoffel symbols of the second kind) are:
\begin{equation}
\Gamma^\rho_{\mu\nu}(v) = \frac{1}{2} \sum_\sigma g^{\rho\sigma}(v) \Bigl(\partial_\mu g_{\nu\sigma}(v) + \partial_\nu g_{\mu\sigma}(v) - \partial_\sigma g_{\mu\nu}(v)\Bigr)
\end{equation}
where $g^{\rho\sigma}$ is the inverse metric (diagonal in this construction).
\end{definition}

\begin{lemma}[Flat Spacetime Has Zero Connection {\normalfont (EinsteinEquations4D.v)}, \textbf{(S)}]
For uniform mass distribution:
\begin{equation}
\forall w,\; \text{mass}(w) = m \quad \implies \quad \Gamma^\rho_{\mu\nu}(v) = 0
\end{equation}
\end{lemma}

\begin{proof}
Uniform mass $\implies$ position-independent metric $\implies$ $\partial_\mu g_{\nu\sigma} = 0$ $\implies$ $\Gamma^\rho_{\mu\nu} = 0$.
\end{proof}

\subsection{Curvature Tensors}

\begin{definition}[Riemann Curvature Tensor {\normalfont (RiemannTensor4D.v)}]
The Riemann curvature tensor measures the failure of parallel transport:
\begin{align}
R^\rho_{\sigma\mu\nu}(v) &= \partial_\mu \Gamma^\rho_{\nu\sigma}(v) - \partial_\nu \Gamma^\rho_{\mu\sigma}(v) \nonumber \\
&\quad + \sum_\lambda \Bigl(\Gamma^\rho_{\mu\lambda}(v)\Gamma^\lambda_{\nu\sigma}(v) - \Gamma^\rho_{\nu\lambda}(v)\Gamma^\lambda_{\mu\sigma}(v)\Bigr)
\end{align}
\end{definition}

\begin{definition}[Ricci Tensor {\normalfont (RiemannTensor4D.v)}]
Contraction of the Riemann tensor:
\begin{equation}
R_{\mu\nu}(v) = \sum_\rho R^\rho_{\mu\rho\nu}(v)
\end{equation}
\end{definition}

\begin{definition}[Ricci Scalar {\normalfont (RiemannTensor4D.v)}]
Complete contraction:
\begin{equation}
R(v) = \sum_{\mu,\nu} g^{\mu\nu}(v) R_{\mu\nu}(v)
\end{equation}
\end{definition}

\begin{definition}[Einstein Tensor {\normalfont (RiemannTensor4D.v)}]
The Einstein tensor encodes spacetime curvature:
\begin{equation}
G_{\mu\nu}(v) = R_{\mu\nu}(v) - \frac{1}{2} g_{\mu\nu}(v) R(v)
\end{equation}
\end{definition}

\begin{lemma}[Flat Spacetime Has Zero Curvature {\normalfont (EinsteinEquations4D.v)}, \textbf{(S)}]
For uniform mass:
\begin{equation}
\forall w,\; \text{mass}(w) = m \quad \implies \quad R^\rho_{\sigma\mu\nu}(v) = 0 \quad \implies \quad G_{\mu\nu}(v) = 0
\end{equation}
\end{lemma}

\subsection{Stress-Energy Tensor}

\begin{remark}[The components below are named by analogy, not by physics]
\label{rem:stress_energy_misnomer}
The three definitions that follow assign the letters $T_{00}$, $T_{0i}$,
$T_{ij}$ to expressions in \texttt{module\_\allowbreak{}structural\_\allowbreak{}mass}. The names are
borrowed from the relativistic stress-energy tensor; the objects are not those
quantities, and three specific mismatches should be stated plainly rather than
left for a reader to discover:

\begin{enumerate}
  \item \textbf{$T_{00}$ is not an energy density.} It is set equal to
    $\mathrm{mass}(v)$, a dimensionless integer attached to a vertex. An energy
    density is energy per unit volume; there is no volume element anywhere in
    this construction, and no conversion constant. Compare
    Section~\ref{sec:dimensional_gap}: a dimensionless integer ledger cannot
    produce a dimensional constant, which is exactly the obstruction here.
  \item \textbf{$T_{0i}$ is not a momentum density.} Momentum density is
    $\rho v^i$, mass density times velocity. The definition below is
    $\partial_i T_{00}$, the spatial \emph{gradient} of the energy-density
    slot. That is a different object with different index structure and
    different dimensions; a gradient of a density is not a flux. There is no
    velocity field in this model for it to be built from.
  \item \textbf{$T_{ij} = \delta_{ij}\cdot\mathrm{mass}$ silently fixes an
    equation of state.} An isotropic spatial block with pressure equal to
    energy density is $p = \rho$: a \emph{stiff fluid}, the maximally rigid
    equation of state consistent with causality in the relativistic setting.
    Nothing in the model motivates that choice; it is a definitional
    convenience, and it should not be read as a derived result about the
    matter content.
\end{enumerate}

None of this affects the vacuum theorem, where every component is $0$ because
$\mathrm{mass} \equiv 0$. That is also why these mismatches have never
mattered computationally. They matter for how the section is read.
\end{remark}

\begin{definition}[$T_{00}$ slot: vertex mass {\normalfont (EinsteinEquations4D.v)}]
The $00$ component is set to the computational mass. Per
Remark~\ref{rem:stress_energy_misnomer} this is not an energy density:
\begin{equation}
T_{00}(v) = \text{mass}(v)
\end{equation}
\end{definition}

\begin{definition}[$T_{0i}$ slot: spatial gradient of $T_{00}$ {\normalfont (EinsteinEquations4D.v)}]
The mixed components are set to the spatial derivative of the $00$ slot. Per
Remark~\ref{rem:stress_energy_misnomer} this is a gradient, \emph{not} a
momentum density:
\begin{equation}
T_{0i}(v) = T_{i0}(v) = \partial_i T_{00}(v)
\end{equation}
\end{definition}

\begin{definition}[$T_{ij}$ slot: isotropic, $p = \rho$ {\normalfont (EinsteinEquations4D.v)}]
The spatial block is isotropic with diagonal entries equal to the $00$ slot.
Per Remark~\ref{rem:stress_energy_misnomer}, setting the pressure equal to the
energy-density slot is the stiff-fluid equation of state $p = \rho$; it is
imposed here by definition, not derived:
\begin{equation}
T_{ij}(v) = \begin{cases}
\text{mass}(v) & \text{if } i = j \\
0 & \text{if } i \neq j
\end{cases}
\end{equation}
\end{definition}

\begin{definition}[Stress-Energy Tensor {\normalfont (EinsteinEquations4D.v)}]
Combining all components:
\begin{equation}
T_{\mu\nu}(s, \mathcal{K}, v) = \begin{cases}
T_{00}(v) & \text{if } \mu = \nu = 0 \\
T_{0i}(v) & \text{if } \mu = 0, \nu = i \text{ or } \mu = i, \nu = 0 \\
T_{ij}(v) & \text{if } \mu = i, \nu = j \\
0 & \text{otherwise}
\end{cases}
\end{equation}
\end{definition}

\subsection{Einstein's Field Equations}

\begin{definition}[Gravitational Constant {\normalfont (EinsteinEquations4D.v)}]
In computational units where the fundamental scale is set by $\mu$-costs:
\begin{equation}
G = \frac{1}{8\pi}
\end{equation}
This makes $8\pi G = 1$, normalizing the coupling in Einstein's equations.
\end{definition}

\begin{lemma}[Curvature from $\mu$-Gradients {\normalfont (EinsteinEquations4D.v)}, \textbf{(S)}]
\label{lem:curvature_mu}
For uniform mass distribution:
\begin{equation}
\forall w,\; \text{mass}(w) = m \quad \implies \quad G_{\mu\nu}(v) = 0
\end{equation}
\end{lemma}

\begin{proof}
The proof chain is:
\begin{enumerate}
    \item Uniform mass $\implies$ position-independent metric
    \item Position-independent metric $\implies$ zero discrete derivatives $\partial_\mu g_{\nu\sigma} = 0$
    \item Zero metric derivatives $\implies$ zero Christoffel symbols $\Gamma^\rho_{\mu\nu} = 0$
    \item Zero Christoffel $\implies$ zero Riemann tensor $R^\rho_{\sigma\mu\nu} = 0$
    \item Zero Riemann $\implies$ zero Ricci tensor and scalar
    \item Zero Ricci $\implies$ zero Einstein tensor $G_{\mu\nu} = 0$
\end{enumerate}
All steps are proven constructively in \texttt{EinsteinEquations4D.v}.
\end{proof}

\begin{lemma}[Vacuum Implies Zero Stress-Energy {\normalfont (EinsteinEquations4D.v)}, \textbf{(S)}]
\label{lem:stress_conservation}
For vacuum states (zero mass everywhere):
\begin{equation}
\forall w,\; \text{mass}(w) = 0 \quad \implies \quad T_{\mu\nu}(v) = 0
\end{equation}
\end{lemma}

\begin{proof}
Direct computation:
\begin{itemize}
    \item $T_{00} = \text{mass} = 0$
    \item $T_{0i} = \partial_i(\text{mass}) = \partial_i(0) = 0$
    \item $T_{ij} = \delta_{ij} \cdot \text{mass} = \delta_{ij} \cdot 0 = 0$
\end{itemize}
\end{proof}

\begin{theorem}[Vacuum Case: Both Sides Are Zero {\normalfont (\texttt{coq/\allowbreak{}kernel/\allowbreak{}curvature/\allowbreak{}EinsteinEquations4D.v})}, \textbf{(R)}]
\label{thm:einstein_main}
For any VM state $s$, 4D simplicial complex $\mathcal{K}$, spacetime indices $\mu, \nu$, and vertex $v$:
\begin{equation}
\forall w,\; \text{mass}(w) = 0 \quad \implies \quad G_{\mu\nu}(s,\mathcal{K},v) = 8\pi G \, T_{\mu\nu}(s,\mathcal{K},v)
\end{equation}
\textbf{Honest status: this is $0 = 0$.} Both sides equal zero in vacuum because all masses vanish. The gravitational constant $G := 1/(8\pi)$ is a definition (unit choice), not a derived quantity. This is \textbf{not} a derivation of Einstein's field equations. It is a consistency check that the discrete geometric objects reduce correctly when mass is everywhere zero. The non-vacuum case is not proven. This theorem does not appear in the primary claim ledger of the project.
\end{theorem}

\begin{proof}
\textbf{Step 1}: Expand the right-hand side using $G = 1/(8\pi)$:
\begin{equation}
8\pi G \, T_{\mu\nu} = 8\pi \cdot \frac{1}{8\pi} \cdot T_{\mu\nu} = T_{\mu\nu}
\end{equation}

\textbf{Step 2}: Apply Lemma~\ref{lem:curvature_mu} (curvature from $\mu$-gradients):

Vacuum (mass $= 0$ everywhere) is a special case of uniform mass distribution, so:
\begin{equation}
G_{\mu\nu}(v) = 0
\end{equation}

\textbf{Step 3}: Apply Lemma~\ref{lem:stress_conservation} (stress-energy conservation):

For vacuum:
\begin{equation}
T_{\mu\nu}(v) = 0
\end{equation}

\textbf{Conclusion}: Both sides equal zero:
\begin{equation}
0 = 0 \quad \checkmark
\end{equation}

The complete proof is formalized in \texttt{coq/\allowbreak{}kernel/\allowbreak{}curvature/\allowbreak{}EinsteinEquations4D.v} with:
\begin{itemize}
    \item No project-local axioms
    \item Zero admitted lemmas
    \item Zero classical logic imports
    \item Verified by automated proof auditing (Inquisitor: 0 HIGH findings)
\end{itemize}
\end{proof}

\begin{remark}[Vacuum vs. Non-Vacuum]
Theorem~\ref{thm:einstein_main} requires the vacuum hypothesis (mass $= 0$ everywhere). For non-vacuum states with arbitrary mass distributions, the equation $G_{\mu\nu} = 8\pi G T_{\mu\nu}$ only holds when the distribution satisfies discrete Bianchi identities (conservation laws). This is analogous to classical general relativity, where Einstein's equations are \emph{field equations} that constrain which configurations are physically possible, not universal identities satisfied by all states.

The vacuum construction is an internal consistency result. It does not establish that physical spacetime curvature emerges from computation. A non-vacuum physical derivation requires additional geometric and physical premises with independently justified applicability.
\end{remark}

\begin{remark}[Epistemological Status: \textbf{(R)} Consistency Relation]
This is a \textbf{(R)} result: the vacuum theorem reduces to $0 = 0$ (both sides vanish when all masses are zero; the gravitational constant is a definition). It establishes that the discrete geometric objects are internally consistent in the zero-mass case. It does \emph{not} establish that spacetime geometry emerges from computation in any physically meaningful sense; that would require the A\_\allowbreak{}QM named hypothesis and is an open empirical question. The non-vacuum case (the full Einstein field equations for arbitrary mass distributions) is not proved and is not a claim of this project.
\end{remark}

\begin{remark}
The discrete geometric objects (metric, curvature, Einstein tensor) are defined from computational primitives. In the vacuum case (all masses zero) they all reduce to zero. This is a consistency check, not a derivation of spacetime emergence. The claim that spacetime geometry is an emergent property of computation would require more: a non-trivial non-vacuum case, a connection to physical observations, and the named A\_\allowbreak{}QM bridge. None of those are present here.
\end{remark}

\subsection{Empirical Validation}

The theoretical results are validated numerically:

\begin{theorem}[Numerical Verification {\normalfont (tests/test\_\allowbreak{}*)}, Empirical]
\leavevmode
\begin{enumerate}
    \item \textbf{Flat spacetime} (uniform mass $m=1$):
        \begin{itemize}
            \item $\max|\Gamma^\rho_{\mu\nu}| < 10^{-10}$ $\checkmark$
            \item $\max|R^\rho_{\sigma\mu\nu}| < 10^{-10}$ $\checkmark$
            \item $\max|G_{\mu\nu}| < 10^{-10}$ $\checkmark$
        \end{itemize}
    \item \textbf{Curved spacetime} (non-uniform mass $m \in [0.5, 1.5]$):
        \begin{itemize}
            \item $\max|\Gamma^\rho_{\mu\nu}| \approx 3.0$ $\checkmark$ (curvature detected)
            \item $\max|R^\rho_{\sigma\mu\nu}| > 0$ $\checkmark$
        \end{itemize}
    \item \textbf{Vacuum Einstein equations} (mass $= 0$):
        \begin{itemize}
            \item $|G_{\mu\nu}| < 10^{-10}$ $\checkmark$
            \item $|T_{\mu\nu}| < 10^{-10}$ $\checkmark$
            \item $|G_{\mu\nu} - 8\pi G T_{\mu\nu}| < 10^{-10}$ $\checkmark$
        \end{itemize}
\end{enumerate}
\end{theorem}

These tests confirm that the discrete differential geometry correctly implements the theoretical predictions.

\subsection{Affine Metric-Scaled Operator: Non-Uniform Boundary Results}

The first-neighbor discrete operator refutes the non-uniform diagonal extension of the EFE. A signed affine metric-scaled redesign closes the gap at vertex 1 and fully characterizes the outer vertices.

\begin{theorem}[Affine Off-Diagonal Ricci Zero {\normalfont (kernel/AffineEFEClosure.v)}, \textbf{(S)}]
\texttt{affine\_\allowbreak{}off\_\allowbreak{}diagonal\_\allowbreak{}ricci\_\allowbreak{}zero}: Under the affine metric-scaled symmetric operator, with the non-uniform isotropic boundary 4-simplex metric, off-diagonal Ricci vanishes at vertex 1: $R_{\mu\nu}(v_1) = 0$ for $\mu \neq \nu$. Zero Admitted. Zero Section Variables.
\end{theorem}

\begin{theorem}[Full Tensor EFE at Vertex 1 {\normalfont (kernel/AffineEFEClosure.v)}, \textbf{(S)}]
\texttt{affine\_\allowbreak{}full\_\allowbreak{}tensor\_\allowbreak{}efe}: Under \texttt{module\_\allowbreak{}structural\_\allowbreak{}mass s 1 = 1} (unit structural mass at vertex 1), the full tensor Einstein field equation holds:
\[ G_{\mu\nu}(v_1) = 3 \cdot T_{\mu\nu}(v_1) \]
for all $\mu, \nu \in \{0,1,2,3\}$. Zero Admitted.
\end{theorem}

\begin{theorem}[Outer Vertex Characterization {\normalfont (kernel/AffineEFEClosure.v)}, \textbf{(S)}]
At the outer vertices $v \in \{0,2,3,4\}$ of the boundary 4-simplex, the affine operator produces:
\begin{itemize}
    \item \texttt{affine\_\allowbreak{}ricci\_\allowbreak{}offdiag\_\allowbreak{}outer}: Off-diagonal Ricci $= 3/32$ (non-zero).
    \item \texttt{affine\_\allowbreak{}ricci\_\allowbreak{}diag\_\allowbreak{}outer}: Diagonal Ricci $= -45/32$.
    \item \texttt{affine\_\allowbreak{}ricci\_\allowbreak{}scalar\_\allowbreak{}outer}: Ricci scalar $= -45/8$.
    \item \texttt{affine\_\allowbreak{}einstein\_\allowbreak{}outer}: Einstein tensor $= 45/32$ (diagonal) and $3/32$ (off-diagonal).
    \item \texttt{affine\_\allowbreak{}efe\_\allowbreak{}fails\_\allowbreak{}outer\_\allowbreak{}offdiag}: The diagonal EFE form $G = 3T$ fails at outer vertices because $G$ has non-zero off-diagonal components.
\end{itemize}
The affine factor $8 - 5g_{00}$ is specifically tuned to the center vertex geometry; at outer vertices it produces a non-trivial off-diagonal Einstein tensor. Zero Admitted.
\end{theorem}

\begin{remark}[What Remains Open]
The affine operator closes vertex 1 and characterizes the outer vertices. The remaining open question is whether the affine factor $8 - 5g_{00}$ is the correct physical discretization. That requires connecting the discrete operator to a physical substrate, not a Coq question, an empirical one.
\end{remark}

\chapter{Information Causality}

\begin{theorem}[IC--$\mu$ Equivalence {\normalfont (InformationCausality.v)}, \textbf{(R)}]
The file defines two record types (\texttt{ICScenario} and \texttt{MuScenario}) with opaque \texttt{Prop} fields, then proves they are equivalent \emph{given} a hypothesis (\texttt{ic\_\allowbreak{}mu\_\allowbreak{}equivalent}) that literally states the equivalence. The proof is \texttt{tauto}. No mutual information, probability distributions, or CHSH values appear anywhere in the formalization.
\end{theorem}

\begin{remark}
The intended statement, that Bob's total information about Alice's data is bounded by the channel's $\mu$-capacity, is a reasonable physical claim but is \textbf{not formalized} in the Coq file. The current file is a placeholder that assumes the conclusion as a hypothesis and re-derives it. A future version would need to define entropy, mutual information, and prove the bound from the VM semantics.
\end{remark}

\chapter{Thermodynamic Bridge}

\section{Landauer's Principle}

\begin{theorem}[Landauer Derived {\normalfont (LandauerDerived.v)}]
\leavevmode
\begin{enumerate}
    \item \texttt{erasure\_\allowbreak{}irreversible}: Erasing $\ge 1$ bit implies fan-in $= 2^k > 1$ (arithmetic: $2^k > 1$ for $k \ge 1$).
    \item \texttt{erasure\_\allowbreak{}decreases\_\allowbreak{}entropy}: Erasure decreases system entropy ($\Delta S < 0$), this is just $\mathrm{output} < \mathrm{input} \Rightarrow \mathrm{output} - \mathrm{input} < 0$.
    \item \texttt{landauer\_\allowbreak{}information\_\allowbreak{}bound}: For any physical erasure, environment entropy increase $\ge$ bits erased.
    \item \texttt{erasure\_\allowbreak{}additive}: Sequential erasures compose additively.
\end{enumerate}

\textbf{Caveat on item~3}: The \texttt{PhysicalErasure} record has \texttt{second\_\allowbreak{}law\_\allowbreak{}satisfied} as a \textbf{field}, not a derived property. The theorem \texttt{landauer\_\allowbreak{}information\_\allowbreak{}bound} simply extracts this field, making the bound an \emph{input assumption} rather than a derivation. The physical energy bound $Q \ge k_B T \ln 2 \cdot n$ is \textbf{explicitly not proven} (the file itself notes the conversion factor is not formalized). Also, the \texttt{mu} in this file is a standalone synonym for \texttt{bits\_\allowbreak{}erased}, it has no formal connection to \texttt{vm\_\allowbreak{}mu} in \texttt{VMState.v}.
\end{theorem}

\section{A2 from Landauer (Factored Form)}
\label{sec:f1}

\begin{theorem}[F1 Strong Form: A2 via Physical Landauer {\normalfont (\texttt{coq/\allowbreak{}kernel/\allowbreak{}frontier/\allowbreak{}F1\_\allowbreak{}StrongForm.v})}, \textbf{(C)}]
\label{thm:f1_strong_form}
Given an abstract dissipation function $\phi : \mathrm{vm\_\allowbreak{}instruction} \to \N$ (interpreted as dissipation per step in Landauer quanta), if
\begin{enumerate}
    \item \textbf{Landauer's principle} holds in dissipation vocabulary, every class-collapsing step on any bool macro-property has $\phi(i) \geq 1$, and
    \item the named bridge \texttt{cost\_\allowbreak{}dissipation\_\allowbreak{}calibrated} holds, $\mathrm{instruction\_\allowbreak{}cost}(i) \geq \phi(i)$ for every $i$,
\end{enumerate}
then on every cert-flip step $\mathrm{instruction\_\allowbreak{}cost}(i) \geq 1$. The Landauer-style premise mentions only class-collapse and dissipation; \texttt{instruction\_\allowbreak{}cost} appears only in the calibration. This is a valid conditional implication, but its two premises are incompatible on the current full VM instruction set, as the following remark explains.
\end{theorem}

\begin{remark}[Applicability failure of the full-ISA physical premises]
\texttt{F1\_\allowbreak{}physical\_\allowbreak{}premises\_\allowbreak{}incompatible} proves that no dissipation function satisfies both hypotheses. The instruction \texttt{JUMP}(1,0) has zero cost and sends every state's program counter to $1$. For the Boolean macro-property ``PC equals $1$,'' it meets \texttt{step\_\allowbreak{}collapses\_\allowbreak{}bool\_\allowbreak{}classes}: a state with PC $0$ enters the true class, and states already in that class stay there. Hypothesis (1) therefore requires dissipation at least $1$, while hypothesis (2) bounds it by $0$.

Thus this theorem does not supply an applicable physical derivation of A2 for the current VM. A2 follows directly from its instruction semantics, without either physical premise. A physical account must justify a domain of operations and macro-properties, reconcile its dissipation model with the cost schedule, and exhibit jointly satisfiable premises. Merely losing a distinction under a selected macro-observation does not itself establish physical information erasure.
\end{remark}

\begin{remark}[Integrality of $\phi$ is a physical assumption, not a typing choice]
\label{rem:phi_integrality}
$\phi$ is declared with codomain $\N$, and the Landauer premise is stated as
$\phi(i) \geq 1$. The integrality is doing real work: it is what makes the
floor come out as ``at least one'' rather than ``at least $k_B T \ln 2$''.
But Landauer's principle gives a \emph{real-valued} lower bound on dissipated
energy per bit erased; nothing in thermodynamics says dissipation arrives in
integer multiples of a quantum. Choosing $\N$ as the codomain therefore
imports a quantisation assumption through the type signature, where it is easy
to miss. It should be read as a named modelling assumption on the same footing
as \texttt{cost\_\allowbreak{}dissipation\_\allowbreak{}calibrated}, not as a neutral formalisation
detail.
\end{remark}

\begin{theorem}[Calibration consistency and retained history, \textbf{(S)}]
For eligible operations $E(i)$ and a schedule $c(i)$, a function $d$ with $E(i)\Rightarrow d(i)\geq1$ and $d(i)\leq c(i)$ exists exactly when all eligible operations have positive scheduled cost. The witness in the converse is $d=c$. A Boolean identity/reset model realizes the conditions with charges zero and one respectively.

For any deterministic step function, the extension
\[
(s,h)\mapsto(\operatorname{step}(s,i),s::h)
\]
is injective for fixed $i$, has a left inverse on its image, and projects to the original transition. In particular, the zero-cost VM jump has an injective history lift. These results establish consistency of an abstract restricted model and failure of macro-collapse as a general erasure test. They do not identify the cost units with heat.
\end{theorem}

\section{The Dimensional Gap: Why Ledgers Cannot Fix Constants}
\label{sec:dimensional_gap}

\begin{theorem}[Dimensional Gap {\normalfont (\texttt{coq/\allowbreak{}kernel/\allowbreak{}thermodynamic/\allowbreak{}DimensionalGapTheorem.v})}, \textbf{(S)}]
\label{thm:dimensional_gap}
Let $\mathrm{State}$ be any type, $L : \mathrm{State} \to \N$ any dimensionless
integer ledger, $\omega : \mathrm{State} \to \N$ any microstate-count function,
and $S_{\mathrm{cand}} : \mathrm{State} \to \mathbb{R}$ any candidate entropy. Suppose
for some $\mathrm{base} \geq 2$ and some state $s$ with $L(s) \geq 1$:
\begin{enumerate}
    \item \textbf{ledger-exponential}: $\omega(s') = \mathrm{base}^{L(s')}$ for all $s'$;
    \item \textbf{ledger-linear}: $S_{\mathrm{cand}}(s') = \alpha \cdot L(s')$ for all $s'$;
    \item \textbf{Boltzmann-form}: $S_{\mathrm{cand}}(s') = k_B \ln \omega(s')$ whenever $\omega(s') > 0$.
\end{enumerate}
Then $\alpha$ is \emph{forced}:
\begin{equation}
\alpha = k_B \cdot \ln(\mathrm{base}).
\end{equation}
\end{theorem}

This is the sharpest honest statement of the limit on deriving physics from
computation in this framework, and it deserves to be read as a \textbf{no-go},
not as a bridge. The theorem quantifies over \emph{any} type, \emph{any}
integer ledger, and \emph{any} exponential microstate count. It makes no
reference to the Thiele machine at all. What it says is that a dimensionless
integer ledger determines the entropy relation only \emph{up to} two externally
supplied scalars: the unit-conversion constant $k_B$ (an empirical fact about
thermodynamics) and the counting base (a structural fact about the substrate's
degrees of freedom). Fix both and $\alpha$ follows; fix neither and the ledger
does not prefer any value. There is a one-parameter family of consistent
$(\alpha, k_B)$ pairs, and choosing among them is precisely the modelling input
that the ledger cannot supply.

Two consequences worth stating explicitly:

\begin{itemize}
  \item The $k_B \ln 2$ appearing in the Thiele instantiation
    (\texttt{SecondLawBoltzmannWall.v}, base $= 2$ for binary cert-flips) is
    \emph{not} derived from the machine. It is what you get after supplying
    $k_B$ and the base from outside.
  \item Any claim of the form ``this computational framework predicts a
    physical constant'' is blocked by Theorem~\ref{thm:dimensional_gap} unless
    it also says which two external scalars it assumed. This applies to the
    present development as much as to any other.
\end{itemize}

Read this section before the conditional physics chapters, not after: it is
the constraint they operate under, and it is what keeps them honest.

\chapter{Self-Reference and the Thiele Manifold}

\section{Self-Reference Escalation}

\begin{theorem}[G\"odel-Style Escalation {\normalfont (SelfReference.v)}]
Any self-referential system $S$ requires a meta-system $\mathrm{Meta}$ with strictly more dimensions:
\[ \mathrm{self\_\allowbreak{}ref}(S) \implies \exists\, \mathrm{Meta}: \dim(\mathrm{Meta}) > \dim(S) \]
The meta-system inherits self-reference, creating an infinite escalation.
\end{theorem}

\section{The Thiele Manifold}

\begin{definition}[Thiele Manifold {\normalfont (ThieleManifold.v)}]
An infinite tower of systems $\{L_n\}_{n \in \N}$ with:
\begin{itemize}
    \item $\dim(L_n) = 4 + n$ (level $n$ has dimension $4+n$).
    \item $\dim(L_{n+1}) > \dim(L_n)$ (strict enrichment).
    \item Each level reasons about the level below.
    \item Self-reference at level $n$ is answered by level $n+1$.
\end{itemize}
\end{definition}

\begin{theorem}[Projection to Spacetime {\normalfont (ThieleManifold.v)}]
The projection $\pi_4$ collapses the tower to dimension 4 (spacetime). For $n > 0$:
\begin{itemize}
    \item $\pi_4$ is lossy: $\dim(L_n) > 4$.
    \item $\mu$-cost of projection $= \dim(L_n) - 4 = n > 0$.
\end{itemize}
Spacetime is the ``shadow'' of the full manifold.
\end{theorem}
\chapter{Three-Layer Observable Agreement}

\begin{remark}[Agreement, not isomorphism]
No isomorphism is constructed in this chapter, and none is claimed: there is no
pair of mutually inverse maps between the Coq, OCaml/Python, and RTL layers, and
no proof that either composite is the identity. What is established is that each
layer's decoder satisfies a shared observable-projection contract, and that the
three projections agree on every tested input: transitivity of equality
through a common projection.

The distinction is not pedantic. Agreement on a projection is compatible with
the layers differing arbitrarily on everything the projection forgets; an
isomorphism would forbid that. Since the projection deliberately forgets
internal representation, agreement is the right property to want here. Isomorphism would be a stronger claim than the construction supports.
\end{remark}

\begin{theorem}[Observable Comparison Contract {\normalfont (ThreeLayerIsomorphism.v)}]
The active Coq theorem packages the \emph{observable-state contract} used for
cross-layer comparison: if the chosen Python and RTL decoders satisfy the same
projection contract as the Coq layer, then the decoded observable states agree
on that projection for a shared trace $\tau$.
\[
\mathrm{decode}_{\mathrm{Python}}(\tau) = \mathrm{decode}_{\mathrm{Coq}}(\tau)
\;\wedge\;
\mathrm{decode}_{\mathrm{RTL}}(\tau) = \mathrm{decode}_{\mathrm{Coq}}(\tau)
\Longrightarrow
\mathrm{decode}_{\mathrm{Python}}(\tau) = \mathrm{decode}_{\mathrm{RTL}}(\tau)
\]
This is a contract theorem over explicit observable projections, not an
unconditional claim that every current Python and RTL artifact is universally
equal to Coq in all CI environments.
\end{theorem}

\section{RTL Coverage: Zero Gaps}

\begin{theorem}[RTL Gap Registry Empty {\normalfont (kami\_\allowbreak{}hw/RTLGapRegistry.v)}, \textbf{(S)}]
\texttt{rtl\_\allowbreak{}gap\_\allowbreak{}count}: The RTL gap registry is empty. There are zero unresolved coverage gaps between the Kami hardware specification and the Coq kernel semantics:
\[ \texttt{rtl\_\allowbreak{}gap\_\allowbreak{}count} = 0 \]
\texttt{rtl\_\allowbreak{}coverage\_\allowbreak{}partition}: The 47 synth-realized opcodes (of the 51-opcode ISA) partition as:
\[ 37 + 10 + 0 = 47. \]
The counts denote unconditional cases, cases under the joint representation invariant, and registered gaps respectively. The joint invariant includes \texttt{coupling\_\allowbreak{}desc\_\allowbreak{}safe} as well as \texttt{coupling\_\allowbreak{}wf} and \texttt{morph\_\allowbreak{}table\_\allowbreak{}wf}.
This is a Coq theorem (\texttt{rtl\_\allowbreak{}coverage\_\allowbreak{}partition} in \texttt{coq/\allowbreak{}kami\_\allowbreak{}hw/\allowbreak{}RTLGapRegistry.v}), not a comment. The partition equation is machine-checked. Zero Admitted.
\end{theorem}

\begin{remark}
The official partition (\texttt{rtl\_\allowbreak{}coverage\_\allowbreak{}partition} in \texttt{coq/\allowbreak{}kami\_\allowbreak{}hw/\allowbreak{}RTLGapRegistry.v}) is $37 + 10 + 0 = 47$. The 37 unconditional proofs are: 31 from the SupportedOpcode group, including \texttt{CHSH\_\allowbreak{}LASSERT} (\texttt{coq/\allowbreak{}kami\_\allowbreak{}hw/\allowbreak{}EmbedStep.v}: \texttt{embed\_\allowbreak{}step\_\allowbreak{}compute}) plus CALL, RET, CHSH\_\allowbreak{}TRIAL (\texttt{coq/\allowbreak{}kami\_\allowbreak{}hw/\allowbreak{}EmbedStep\_\allowbreak{}WF.v}) and TENSOR\_\allowbreak{}SET, TENSOR\_\allowbreak{}GET, LASSERT (\texttt{coq/\allowbreak{}kami\_\allowbreak{}hw/\allowbreak{}GraphReconstructionBridge.v}). The remaining 10 (PNEW, PSPLIT, PMERGE, MORPH, MORPH\_\allowbreak{}ID, MORPH\_\allowbreak{}DELETE, MORPH\_\allowbreak{}ASSERT, MORPH\_\allowbreak{}GET, COMPOSE, MORPH\_\allowbreak{}TENSOR) carry Qed proofs under the joint structural invariant \texttt{morph\_\allowbreak{}table\_\allowbreak{}wf} $\wedge$ \texttt{coupling\_\allowbreak{}wf} $\wedge$ \texttt{coupling\_\allowbreak{}desc\_\allowbreak{}safe}. Each component is preserved by every \texttt{kami\_\allowbreak{}step}: \texttt{morph\_\allowbreak{}table\_\allowbreak{}wf\_\allowbreak{}kami\_\allowbreak{}step\_\allowbreak{}preserved}, \texttt{coupling\_\allowbreak{}desc\_\allowbreak{}safe\_\allowbreak{}kami\_\allowbreak{}step\_\allowbreak{}preserved}, and \texttt{coupling\_\allowbreak{}wf\_\allowbreak{}kami\_\allowbreak{}step\_\allowbreak{}preserved}. The last takes \texttt{coupling\_\allowbreak{}desc\_\allowbreak{}safe} as a side hypothesis (in addition to \texttt{coupling\_\allowbreak{}wf}), which is why the conjunction is the actual inductive invariant. So the precondition is discharged for any state reachable by machine execution. PNEW/PSPLIT/PMERGE additionally need \texttt{pt\_\allowbreak{}well\_\allowbreak{}formed} and arithmetic side conditions that are likewise inductive.
\end{remark}

\section{Named Trust Boundaries}

The project has exactly two named trust boundaries. They are not Admitted proofs. They are explicit, auditable limits of the formal claim.

\begin{description}
\item[\texttt{bsc\_\allowbreak{}kami\_\allowbreak{}compilation\_\allowbreak{}trusted} (Layer 2, Verilog semantics gap).] The Bluespec SystemVerilog compiler generates Verilog from the Kami hardware specification. The correctness of the BSC compilation step, that the generated Verilog implements exactly the Kami spec, is trusted, not proved in Coq. This is documented in \texttt{coq/\allowbreak{}kernel/\allowbreak{}hardware\_\allowbreak{}bridge/\allowbreak{}VerilogRTLCorrespondence.v} (Section~2b, ``BSC Compilation Trust Boundary'') and the matching \texttt{Definition bsc\_\allowbreak{}kami\_\allowbreak{}compilation\_\allowbreak{}trusted}. The Section Variable \texttt{rtl\_\allowbreak{}step\_\allowbreak{}correct} states that the synthesised RTL implements the Kami step function; this is the residual hypothesis that every theorem in the section quantifies over. The boundary is empirically backed by the cosim suite (\texttt{tests/\allowbreak{}test\_\allowbreak{}verilog\_\allowbreak{}cosim.py}, 38 tests, including \texttt{test\_\allowbreak{}all\_\allowbreak{}47\_\allowbreak{}opcodes\_\allowbreak{}in\_\allowbreak{}cosim} which checks that every opcode appears) and the parametrized cross-layer fuzz suites (\texttt{tests/\allowbreak{}test\_\allowbreak{}fuzz\_\allowbreak{}random\_\allowbreak{}programs.py} and \texttt{tests/\allowbreak{}test\_\allowbreak{}cross\_\allowbreak{}layer\_\allowbreak{}adversarial\_\allowbreak{}fuzz.py}).

\item[\texttt{A\_\allowbreak{}QM} (physical quantum axiom).] Whether a physical Thiele Machine can produce CHSH-violating statistics by exploiting quantum entanglement in hardware is an open empirical question. The \texttt{column\_\allowbreak{}contractive\_\allowbreak{}iff\_\allowbreak{}quantum\_\allowbreak{}realizable} biconditional (Theorem~\ref{thm:thiele_honest_iff_qr}) requires no physical axiom, it is pure algebra. The physical quantum axiom \texttt{A\_\allowbreak{}QM} (declared as a \texttt{Context} parameter in \texttt{coq/\allowbreak{}PhysicsConditionalClosure.v}) is the separate named hypothesis that quantum CHSH experiments produce PSD NPA matrices. It does not appear in the kernel-tier proof of the biconditional; it is the bridge used in conditional theorems such as \texttt{master\_\allowbreak{}tsirelson\_\allowbreak{}conditional}.
\end{description}
\chapter{Categorical Structure}
\label{ch:categorical}

\section{The Morphism Layer}

The active implementation adds a concrete category directly to the machine state. This is not an abstract analogy, it is an executable layer of 7 opcodes (0x27--0x2D) with formally verified laws.

\subsection{Objects and Morphisms}

The partition graph $G$ carries two maps:
\begin{itemize}
    \item $\texttt{pg\_\allowbreak{}modules} : \mathrm{ModuleID} \rightharpoonup \mathrm{ModuleState}$, objects (partition modules)
    \item $\texttt{pg\_\allowbreak{}morphisms} : \mathrm{MorphismID} \rightharpoonup \mathrm{MorphismState}$, arrows
\end{itemize}

A $\mathrm{MorphismState}$ bundles: source module, target module, coupling data (list of $(nat \times nat)$ pairs plus a label), and an identity flag.

\subsection{The Seven Categorical Opcodes}

\begin{itemize}
    \item $\texttt{MORPH}(dst, src, tgt, idx, \Delta\mu)$: create morphism $src \to tgt$, write morphism ID to $\texttt{regs}[dst]$. Costs $\Delta\mu$.
    \item $\texttt{COMPOSE}(dst, f, g, \Delta\mu)$: compose morphisms $f ; g$ when $f.\mathrm{target} = g.\mathrm{source}$, write new ID to $\texttt{regs}[dst]$. Costs $\Delta\mu$. Errors if endpoints mismatch.
    \item $\texttt{MORPH\_\allowbreak{}ID}(dst, m, \Delta\mu)$: create identity morphism $\mathrm{id}_m$ for module $m$ (diagonal relation). Costs $\Delta\mu$.
    \item $\texttt{MORPH\_\allowbreak{}DELETE}(f, \Delta\mu)$: remove morphism $f$. Cascades on module deletion via \texttt{graph\_\allowbreak{}cascade\_\allowbreak{}delete\_\allowbreak{}morphisms}.
    \item $\texttt{MORPH\_\allowbreak{}ASSERT}(f, \mathit{prop}, \mathit{cert}, \Delta\mu)$: cert-setter. Costs $S(\Delta\mu) \geq 1$ unconditionally, even for failed attempts (morphism does not exist). This is the No Free Insight law applied to relational claims.
    \item $\texttt{MORPH\_\allowbreak{}TENSOR}(dst, f, g, \Delta\mu)$: tensor product $f \otimes g$. Requires that the union modules $\mathrm{src}(f) \cup \mathrm{src}(g)$ and $\mathrm{tgt}(f) \cup \mathrm{tgt}(g)$ already exist in the graph.
    \item $\texttt{MORPH\_\allowbreak{}GET}(dst, f, \mathit{sel}, \Delta\mu)$: read field from morphism $f$: selector 0 = source, 1 = target, 2 = coupling length, 3 = is\_\allowbreak{}identity.
\end{itemize}

\subsection{Category Laws}

\begin{theorem}[Relational Composition is Associative {\normalfont (kernel/CategoryLaws.v)}, \textbf{(S)}]
\label{thm:compose_assoc}
For any three morphisms $f : A \to B$, $g : B \to C$, $h : C \to D$ with composable endpoints:
\[ (f ; g) ; h = f ; (g ; h) \]
in the relational sense (coupling pairs of the composite). Zero Admitted.
\end{theorem}

\begin{theorem}[Graph Composition Satisfies Laws {\normalfont (kernel/CategoryBridge.v)}, \textbf{(S)}]
\label{thm:graph_compose_assoc}
The operation $\texttt{graph\_\allowbreak{}compose\_\allowbreak{}morphisms}$ satisfies associativity and unitality. Additionally:
\[ \texttt{categorical\_\allowbreak{}extension\_\allowbreak{}nofi\_\allowbreak{}consistent} \]
MORPH and COMPOSE can carry $\Delta\mu = 0$ and are not in the mandatory-cost class. The NoFI policy enters through $\texttt{MORPH\_\allowbreak{}ASSERT}$: asserting that a built morphism chain exists and activating the cert channel always costs $S(\Delta\mu) \geq 1$. Building the chain is free; certifying it never is.
\end{theorem}

\begin{theorem}[Tensor Bifunctoriality {\normalfont (kernel/CategoryMonoidal.v)}, \textbf{(S)}]
\label{thm:tensor_bifunctor}
$\texttt{graph\_\allowbreak{}tensor\_\allowbreak{}morphisms}$ is a bifunctor: it satisfies the interchange law
\[ (f \otimes g) ; (h \otimes k) = (f ; h) \otimes (g ; k) \]
for composable, disjoint pairs.
\end{theorem}

\subsection{The Categorical Separation Theorem}

\begin{theorem}[Categorical Separation {\normalfont (kernel/PartitionSeparation.v, §10)}, \textbf{(S)}]
\label{thm:categorical_separation}
There exist states $s_1, s_2$ such that:
\[
s_1.\texttt{vm\_\allowbreak{}regs} = s_2.\texttt{vm\_\allowbreak{}regs} \;\wedge\;
s_1.\texttt{vm\_\allowbreak{}mem}  = s_2.\texttt{vm\_\allowbreak{}mem}  \;\wedge\;
s_1.\texttt{vm\_\allowbreak{}mu}   = s_2.\texttt{vm\_\allowbreak{}mu}   \;\wedge\;
s_1.\texttt{vm\_\allowbreak{}err}  = s_2.\texttt{vm\_\allowbreak{}err}
\]
but $s_1.\texttt{vm\_\allowbreak{}graph}.\texttt{pg\_\allowbreak{}morphisms} \neq s_2.\texttt{vm\_\allowbreak{}graph}.\texttt{pg\_\allowbreak{}morphisms}$.
\end{theorem}

\begin{proof}
Construct $s_1$ with $\texttt{pg\_\allowbreak{}morphisms} = [(0, ms)]$ for any morphism state $ms$, and $s_2$ with $\texttt{pg\_\allowbreak{}morphisms} = []$. Arrange identical $\texttt{vm\_\allowbreak{}regs}$, $\texttt{vm\_\allowbreak{}mem}$, $\texttt{vm\_\allowbreak{}mu}$, $\texttt{vm\_\allowbreak{}err}$ by choosing costs and instruction sequences accordingly. Distinction by $\texttt{discriminate}$. \qed
\end{proof}

\begin{remark}[Significance]
A classical machine (exposing only registers, memory, and a cost counter) cannot distinguish $s_1$ from $s_2$. The Thiele Machine distinguishes them in one probe instruction: $\texttt{MORPH\_\allowbreak{}DELETE}\ 0$ succeeds on $s_1$ and errors on $s_2$. The categorical layer records \emph{how} knowledge was derived, not just \emph{what} the final values are.

This is the machine's claim to being a \emph{knowledge auditor}: any party claiming to have verified a structural relation must carry a $\mu$-receipt $\geq$ the minimum construction cost. The NoFI theorem supplies a cost floor under its trace hypotheses. It does not authenticate an externally supplied receipt or prove arbitrary property text true.
\end{remark}
\chapter{Falsifiable Predictions}

The honest falsifiable predictions are the structural theorems that compile or fail. The claims below identify specific mathematical conditions that would constitute counterexamples to the machine's core theorems.

\section{Falsify No Free Insight}

Find a Thiele Machine program trace that starts with \texttt{vm\_\allowbreak{}certified = false} and $\mu = 0$, ends with \texttt{vm\_\allowbreak{}certified = true}, and ends with $\mu = 0$. The Coq proof of \texttt{certification\_\allowbreak{}requires\_\allowbreak{}positive\_\allowbreak{}mu} would fail.

\section{Falsify the Classical CHSH Bound}

Exhibit a locally factorizable strategy achieving $|S| > 2$. This contradicts \texttt{chsh\_\allowbreak{}stat\_\allowbreak{}violation\_\allowbreak{}not\_\allowbreak{}local} and also contradicts known mathematics.

\section{Falsify the Biconditional}

Exhibit $E_{00}, E_{01}, E_{10}, E_{11}$ satisfying the three column-contractivity conditions but with a non-PSD NPA moment matrix. The theorem \texttt{column\_\allowbreak{}contractive\_\allowbreak{}iff\_\allowbreak{}quantum\_\allowbreak{}realizable} says no such example exists.

\section{Falsify the $\mu$-Hierarchy}

Find $k \geq 1$ and a trace reaching \texttt{vm\_\allowbreak{}certified = true} and $\texttt{vm\_\allowbreak{}mu} \geq k$ with total trace cost $< k$. This contradicts \texttt{mu\_\allowbreak{}hierarchy\_\allowbreak{}theorem}.

\section{What I Don't Know}

\begin{itemize}
\item \textbf{Physical interpretation of $\mu$}: Whether the $\mu$ ledger corresponds to thermodynamic entropy in a physical Thiele Machine is an open empirical question. The formal development is consistent with Landauer's principle. The calibration hypothesis is named (\texttt{mu\_\allowbreak{}landauer\_\allowbreak{}unruh\_\allowbreak{}calibrated}) and not proved in Coq.
\item \textbf{Quantum mechanics connection}: The mathematical equivalence between slice coherence and zero-marginal NPA realizability is a kernel-tier theorem. Whether a physical Thiele Machine can produce CHSH-violating statistics by exploiting quantum entanglement in hardware is an open empirical question outside the scope of this project.
\item \textbf{Complexity implications}: Whether $\mu$P $\neq$ $\mu$NP is an open question. The factored-search witness shows a concrete structural advantage but is not a class separation.
\item \textbf{Physical constants}: No physical constant is derived from first principles. There is no derivation of Planck's constant $h$; any relation tying an internal time-scale $\tau_\mu$ to $h$ would be circular by construction, defining $\tau_\mu$ in terms of $h$ and checking consistency rather than predicting $h$.
\item \textbf{The two $+1$s} (OP-Plus-One): the \texttt{+1} in \texttt{triangle\_\allowbreak{}angle}'s denominator (\texttt{PI · d\_\allowbreak{}bc /\allowbreak{} (1 + d\_\allowbreak{}ab + d\_\allowbreak{}ac + d\_\allowbreak{}bc)}, \texttt{coq/\allowbreak{}kernel/\allowbreak{}curvature/\allowbreak{}MuGravity.v}) versus the A2 cost-floor \texttt{+1} on CERTIFY/LASSERT (\texttt{coq/\allowbreak{}kernel/\allowbreak{}foundation/\allowbreak{}VMStep.v}), same axiom, or unrelated regularisation? No kernel-internal theorem unifies them without circularity (\texttt{coq/\allowbreak{}kernel/\allowbreak{}frontier/\allowbreak{}F3\_\allowbreak{}PlusOneStructural.v}); the symbolic correction term decays as $-\pi / (3(3d+1))$ in the equilateral case (\texttt{tools/\allowbreak{}triangle\_\allowbreak{}angle\_\allowbreak{}plus\_\allowbreak{}one\_\allowbreak{}analysis.py}), regularisation-shaped rather than A2-quantum-shaped. Inconclusive. A positive structural identification (Case~A, the geometric \texttt{+1} is A2 in geometric form) would re-frame the calibrated regime's asymptotic-only character as a structural consequence of A2; a positive non-identification (Case~B) would license replacing $S(\dots)$ with $(\dots)$ in the denominator. Either is real categorical content. The kernel was not modified pending closure.
\item \textbf{F3 narrow non-vacuity} (related to OP-Plus-One): a constructible finite \texttt{VMState} satisfying $\forall m, \texttt{calibration\_\allowbreak{}residual}(s, m) = 0$ on a partition graph with at least one actual graph triangle and \texttt{well\_\allowbreak{}formed\_\allowbreak{}triangulated} holding. The symbolic search to $N = 7$ exhausts negatively. Whether the asymptotic-only character of calibration is a structural consequence of OP-Plus-One's Case~A, or just a search-depth artifact, is itself open.
\end{itemize}

%=============================================================================
% APPENDIX
%=============================================================================

\appendix
\chapter{Proof File Index}

The table maps theorem families to their Coq source files in the active build. Observation adequacy, joint event floors, and abstract calibration consistency are in \texttt{coq/\allowbreak{}kernel/\allowbreak{}frontier/\allowbreak{}ObservationPolicy.v}; descent from traces and its concrete targets are in \texttt{coq/\allowbreak{}kernel/\allowbreak{}frontier/\allowbreak{}TraceStateDescent.v}.

\begingroup
\small
\begin{longtable}{@{}>{\raggedright\arraybackslash}p{0.36\textwidth}>{\raggedright\arraybackslash}p{\dimexpr0.64\textwidth-2\tabcolsep\relax}@{}}
\textbf{Result} & \textbf{Coq Source} \\
\hline
\multicolumn{2}{l}{\emph{Cost ledger and No Free Insight}} \\
\hline
No Free Insight (Thm.~\ref{thm:nfi}) & \texttt{coq/\allowbreak{}nofi/\allowbreak{}NoFreeInsight\_\allowbreak{}Theorem.v} \\
$\mu$-Chaitin & \texttt{coq/\allowbreak{}nofi/\allowbreak{}MuChaitinTheory\_\allowbreak{}Theorem.v} \\
$\mu$-Conservation & \texttt{coq/\allowbreak{}kernel/\allowbreak{}foundation/\allowbreak{}MuLedgerConservation.v} \\
Receipt Integrity & \texttt{coq/\allowbreak{}kernel/\allowbreak{}nfi/\allowbreak{}ReceiptIntegrity.v} \\
Noether / Gauge & \texttt{coq/\allowbreak{}kernel/\allowbreak{}curvature/\allowbreak{}KernelNoether.v} \\
Strict Subsumption & \texttt{coq/\allowbreak{}kernel/\allowbreak{}foundation/\allowbreak{}Subsumption.v} \\
Universal NFI (any substrate) & \texttt{coq/\allowbreak{}kernel/\allowbreak{}nfi/\allowbreak{}UniversalCertificationCost.v}: \texttt{universal\_\allowbreak{}nfi\_\allowbreak{}any\_\allowbreak{}substrate} \\
Quantitative NFI ($\ge S(d)$ trace bound) & \texttt{coq/\allowbreak{}kernel/\allowbreak{}mu\_\allowbreak{}calculus/\allowbreak{}QuantitativeNoFI.v} \\
$\mu$-Hierarchy & \texttt{coq/\allowbreak{}kernel/\allowbreak{}mu\_\allowbreak{}calculus/\allowbreak{}MuHierarchyTheorem.v}: \texttt{mu\_\allowbreak{}hierarchy\_\allowbreak{}theorem} \\
Partition Ops Free, Certification Non-Free & \texttt{coq/\allowbreak{}kernel/\allowbreak{}nfi/\allowbreak{}PartitionRefinementNoFI.v} \\
$\mu$-Ledger Necessity (5-way classification) & \texttt{coq/\allowbreak{}NecessityOfMuLedger.v}, \texttt{coq/\allowbreak{}kernel/\allowbreak{}nfi/\allowbreak{}NecessityAbstract.v} \\
Trace Cost for a Fixed Schedule & \texttt{coq/\allowbreak{}NecessityOfMuLedger.v}: \texttt{shadow\_\allowbreak{}mu\_\allowbreak{}is\_\allowbreak{}computation\_\allowbreak{}intrinsic} \\
$\mu$ Cost Universal Across Machines & \texttt{coq/\allowbreak{}NecessityOfMuLedger.v}: \texttt{shadow\_\allowbreak{}mu\_\allowbreak{}delta\_\allowbreak{}universal} \\
$\mu$ Cost Unique Across Accounting Systems & \texttt{coq/\allowbreak{}NecessityOfMuLedger.v}: \texttt{shadow\_\allowbreak{}mu\_\allowbreak{}unique\_\allowbreak{}accounting} \\
Positive-Cost Trace Accumulation & \texttt{coq/\allowbreak{}NecessityOfMuLedger.v}: \texttt{shadow\_\allowbreak{}mu\_\allowbreak{}inevitable} \\
$\mu$-budget predicates + zero-$\mu$ inclusion & \texttt{coq/\allowbreak{}kernel/\allowbreak{}mu\_\allowbreak{}calculus/\allowbreak{}MuComplexity.v} \\
\hline
\multicolumn{2}{l}{\emph{Categorical and shadow-projection layer}} \\
\hline
Categorical Separation / Category Laws & \texttt{coq/\allowbreak{}kernel/\allowbreak{}category/\allowbreak{}CategoryLaws.v}, \texttt{coq/\allowbreak{}kernel/\allowbreak{}category/\allowbreak{}CategoryBridge.v} \\
Shadow Strictly Lossy & \texttt{coq/\allowbreak{}kernel/\allowbreak{}witness/\allowbreak{}ShadowProjection.v}: \texttt{shadow\_\allowbreak{}strictly\_\allowbreak{}lossy} \\
Turing Strictness ($D5$) & \texttt{coq/\allowbreak{}kernel/\allowbreak{}foundation/\allowbreak{}TuringStrictness.v} \\
CHSH Coupling Bridge ($S \le 2$ from local coupling, $S > 2$ rules it out) & \texttt{coq/\allowbreak{}kernel/\allowbreak{}quantum/\allowbreak{}CHSHCouplingBridge.v} \\
Witness Insight Three-Tier Taxonomy & \texttt{coq/\allowbreak{}kernel/\allowbreak{}witness/\allowbreak{}WitnessInsightGeneral.v} \\
\hline
\multicolumn{2}{l}{\emph{Quantum / CHSH algebra (no project-local axioms)}} \\
\hline
Slice-coherent $\iff$ NPA realizable (biconditional) & \texttt{coq/\allowbreak{}kernel/\allowbreak{}quantum/\allowbreak{}QuantumPartitionPSD.v}: \texttt{column\_\allowbreak{}contractive\_\allowbreak{}iff\_\allowbreak{}quantum\_\allowbreak{}realizable} \\
Forward: column-contractive $\Rightarrow$ NPA PSD & \texttt{coq/\allowbreak{}kernel/\allowbreak{}nfi/\allowbreak{}MuLedgerQuantumBridge.v}: \texttt{zero\_\allowbreak{}marginal\_\allowbreak{}npa\_\allowbreak{}column\_\allowbreak{}contractive\_\allowbreak{}implies\_\allowbreak{}psd} \\
Reverse: NPA PSD $\Rightarrow$ column-contractive & \texttt{coq/\allowbreak{}kernel/\allowbreak{}quantum/\allowbreak{}QuantumPartitionPSD.v}: \texttt{npa\_\allowbreak{}psd\_\allowbreak{}implies\_\allowbreak{}column\_\allowbreak{}contractive} \\
Quadratic-form discriminant lemma & \texttt{coq/\allowbreak{}kernel/\allowbreak{}quantum/\allowbreak{}ConstructivePSD.v}: \texttt{quadratic\_\allowbreak{}nonneg\_\allowbreak{}discriminant} \\
Tsirelson from column contractivity ($|S| \le 2\sqrt{2}$) & \texttt{coq/\allowbreak{}kernel/\allowbreak{}nfi/\allowbreak{}MuLedgerQuantumBridge.v}: \texttt{state\_\allowbreak{}column\_\allowbreak{}contractive\_\allowbreak{}implies\_\allowbreak{}tsirelson} (\texttt{tsirelson\_\allowbreak{}from\_\allowbreak{}row\_\allowbreak{}bounds} in \texttt{TsirelsonGeneral.v} is the inner Cauchy--Schwarz step) \\
Tsirelson (algebraic, tight) & \texttt{coq/\allowbreak{}kernel/\allowbreak{}category/\allowbreak{}AlgebraicCoherence.v}: \texttt{tsirelson\_\allowbreak{}bound\_\allowbreak{}tight} \\
Tsirelson (algebraic from coherence) & \texttt{coq/\allowbreak{}kernel/\allowbreak{}quantum/\allowbreak{}TsirelsonFromAlgebra.v} \\
General Algebraic Tsirelson ($S^2 \le 8$) & \texttt{coq/\allowbreak{}kernel/\allowbreak{}category/\allowbreak{}AlgebraicCoherence.v}: \texttt{algebraically\_\allowbreak{}coherent\_\allowbreak{}tsirelson\_\allowbreak{}general} \\
Elliptope realizability (existential completion) & \texttt{coq/\allowbreak{}kernel/\allowbreak{}quantum/\allowbreak{}ElliptopeCompletion.v}: \texttt{elliptope\_\allowbreak{}realizable} (definition); \texttt{deterministic\_\allowbreak{}strategy\_\allowbreak{}elliptope}, \texttt{lhv\_\allowbreak{}mixture\_\allowbreak{}elliptope} (every LHV correlator inside) \\
Tsirelson for the full elliptope ($S^2 \le 8$) & \texttt{coq/\allowbreak{}kernel/\allowbreak{}quantum/\allowbreak{}ElliptopeCompletion.v}: \texttt{elliptope\_\allowbreak{}tsirelson} (inner step \texttt{psd\_\allowbreak{}cauchy\_\allowbreak{}schwarz}) \\
PR box excluded; classical $\subset$ elliptope strict & \texttt{coq/\allowbreak{}kernel/\allowbreak{}quantum/\allowbreak{}ElliptopeCompletion.v}: \texttt{pr\_\allowbreak{}box\_\allowbreak{}not\_\allowbreak{}elliptope}, \texttt{beyond\_\allowbreak{}classical\_\allowbreak{}elliptope} \\
Decidable elliptope gate (soundness) & \texttt{coq/\allowbreak{}kernel/\allowbreak{}quantum/\allowbreak{}ElliptopeGate.v}: \texttt{elliptope\_\allowbreak{}check\_\allowbreak{}full\_\allowbreak{}sound}; boundary acceptance \texttt{gate\_\allowbreak{}accepts\_\allowbreak{}pythagorean\_\allowbreak{}boundary}; universal refusal \texttt{elliptope\_\allowbreak{}full\_\allowbreak{}gate\_\allowbreak{}never\_\allowbreak{}accepts\_\allowbreak{}pr\_\allowbreak{}box} \\
\hline
\multicolumn{2}{l}{\emph{Frontier: pointer-observable criterion}} \\
\hline
Proliferation / unique-pointer schema & \texttt{coq/\allowbreak{}kernel/\allowbreak{}frontier/\allowbreak{}PointerObservable.v}: \texttt{redundantly\_\allowbreak{}proliferating}, \texttt{unique\_\allowbreak{}pointer\_\allowbreak{}among} (definitions); \texttt{toy\_\allowbreak{}cert\_\allowbreak{}unique\_\allowbreak{}pointer} (sanity instance) \\
Five deployed disciplines are unique pointers & \texttt{coq/\allowbreak{}kernel/\allowbreak{}frontier/\allowbreak{}PointerObservableReductions.v}: \texttt{PoS\_\allowbreak{}unique\_\allowbreak{}pointer}, \texttt{Gas\_\allowbreak{}unique\_\allowbreak{}pointer}, \texttt{TEE\_\allowbreak{}unique\_\allowbreak{}pointer}, \texttt{CT\_\allowbreak{}unique\_\allowbreak{}pointer}, \texttt{PCC\_\allowbreak{}unique\_\allowbreak{}pointer}; conjunction \texttt{five\_\allowbreak{}disciplines\_\allowbreak{}are\_\allowbreak{}pointers} (closed under the global context) \\
Adversarial search: forgery-resistant systems that do \emph{not} proliferate & \texttt{coq/\allowbreak{}kernel/\allowbreak{}frontier/\allowbreak{}PointerObservableCounterexamples.v}: \texttt{deniable\_\allowbreak{}authentication\_\allowbreak{}refutes\_\allowbreak{}strong\_\allowbreak{}criterion}, \texttt{mac\_\allowbreak{}refutes\_\allowbreak{}strong\_\allowbreak{}criterion}, \texttt{capability\_\allowbreak{}refutes\_\allowbreak{}strong\_\allowbreak{}criterion}; control \texttt{public\_\allowbreak{}log\_\allowbreak{}confirms} \\
\hline
\multicolumn{2}{l}{\emph{Conditional physics derivations (under named axioms; (C) tier)}} \\
\hline
Born Rule & \texttt{coq/\allowbreak{}kernel/\allowbreak{}quantum/\allowbreak{}BornRule.v} \\
Unitarity & \texttt{coq/\allowbreak{}kernel/\allowbreak{}quantum/\allowbreak{}Unitarity.v} \\
Purification & \texttt{coq/\allowbreak{}kernel/\allowbreak{}quantum/\allowbreak{}Purification.v} \\
No-Cloning & \texttt{coq/\allowbreak{}kernel/\allowbreak{}quantum/\allowbreak{}NoCloning.v} \\
Information Causality & \texttt{coq/\allowbreak{}kernel/\allowbreak{}quantum/\allowbreak{}InformationCausality.v} \\
Landauer derivation & \texttt{coq/\allowbreak{}thermodynamic/\allowbreak{}LandauerDerived.v} \\
Affine Off-Diagonal Ricci Zero (vertex 1) & \texttt{coq/\allowbreak{}kernel/\allowbreak{}curvature/\allowbreak{}AffineEFEClosure.v} \\
Full Tensor EFE at Vertex 1 ($G = 3T$) & \texttt{coq/\allowbreak{}kernel/\allowbreak{}curvature/\allowbreak{}AffineEFEClosure.v} \\
Outer Vertex Characterization (Ricci $= 3/32$) & \texttt{coq/\allowbreak{}kernel/\allowbreak{}curvature/\allowbreak{}AffineEFEClosure.v} \\
Spacetime Emergence (consistency check) & \texttt{coq/\allowbreak{}kernel/\allowbreak{}curvature/\allowbreak{}SpacetimeEmergence.v} \\
$\mu$-Geometry (pseudometric) & \texttt{coq/\allowbreak{}kernel/\allowbreak{}mu\_\allowbreak{}calculus/\allowbreak{}MuGeometry.v} \\
4D Simplicial Complex & \texttt{coq/\allowbreak{}kernel/\allowbreak{}curvature/\allowbreak{}FourDSimplicialComplex.v} \\
Metric from $\mu$-Costs & \texttt{coq/\allowbreak{}kernel/\allowbreak{}curvature/\allowbreak{}MetricFromMuCosts.v} \\
Riemann / Christoffel / Einstein tensors & \texttt{coq/\allowbreak{}kernel/\allowbreak{}curvature/\allowbreak{}RiemannTensor4D.v} \\
Stress-Energy Tensor / Conservation / Vacuum EFE & \texttt{coq/\allowbreak{}kernel/\allowbreak{}curvature/\allowbreak{}EinsteinEquations4D.v} \\
Einstein Equations (general, full diagonal-and-offdiag chain) & \texttt{coq/\allowbreak{}kernel/\allowbreak{}curvature/\allowbreak{}EinsteinEquationsFull.v} \\
\texttt{A\_\allowbreak{}QM} bridge axiom + master Tsirelson conditional & \texttt{coq/\allowbreak{}PhysicsConditionalClosure.v} \\
\hline
\multicolumn{2}{l}{\emph{Self-reference, falsifiability, isomorphism}} \\
\hline
Self-Reference Tower & \texttt{coq/\allowbreak{}self\_\allowbreak{}reference/\allowbreak{}SelfReference.v} \\
Thiele Manifold & \texttt{coq/\allowbreak{}thiele\_\allowbreak{}manifold/\allowbreak{}ThieleManifold.v} \\
Physical Constants (definitions only) & \texttt{coq/\allowbreak{}thiele\_\allowbreak{}manifold/\allowbreak{}PhysicalConstants.v} \\
Falsifiable Prediction record & \texttt{coq/\allowbreak{}kernel/\allowbreak{}aggregators/\allowbreak{}FalsifiablePrediction.v} \\
Three-Layer (Coq $\leftrightarrow$ Python $\leftrightarrow$ RTL) Isomorphism & \texttt{coq/\allowbreak{}kernel/\allowbreak{}hardware\_\allowbreak{}bridge/\allowbreak{}ThreeLayerIsomorphism.v} \\
\hline
\multicolumn{2}{l}{\emph{Hardware / RTL / extraction}} \\
\hline
RTL Gap Count $= 0$ & \texttt{coq/\allowbreak{}kami\_\allowbreak{}hw/\allowbreak{}RTLGapRegistry.v}: \texttt{rtl\_\allowbreak{}gap\_\allowbreak{}count} \\
RTL Coverage Partition ($37 + 10 + 0 = 47$) & \texttt{coq/\allowbreak{}kami\_\allowbreak{}hw/\allowbreak{}RTLGapRegistry.v}: \texttt{rtl\_\allowbreak{}coverage\_\allowbreak{}partition} \\
RTL Trust Boundary (Section Variable) & \texttt{coq/\allowbreak{}kernel/\allowbreak{}hardware\_\allowbreak{}bridge/\allowbreak{}VerilogRTLCorrespondence.v}: \texttt{rtl\_\allowbreak{}step\_\allowbreak{}correct} \\
BSC Compilation Trust Boundary & \texttt{coq/\allowbreak{}kernel/\allowbreak{}hardware\_\allowbreak{}bridge/\allowbreak{}VerilogRTLCorrespondence.v}: \texttt{bsc\_\allowbreak{}kami\_\allowbreak{}compilation\_\allowbreak{}trusted} \\
Coq-Kami refinement (proves \texttt{rtl\_\allowbreak{}step\_\allowbreak{}correct} for COQKAMI) & \texttt{coq/\allowbreak{}kami\_\allowbreak{}hw/\allowbreak{}VerilogSemantics.v} \\
EmbedStep / EmbedStep\_\allowbreak{}WF (per-opcode bisimulation) & \texttt{coq/\allowbreak{}kami\_\allowbreak{}hw/\allowbreak{}EmbedStep.v}, \texttt{coq/\allowbreak{}kami\_\allowbreak{}hw/\allowbreak{}EmbedStep\_\allowbreak{}WF.v} \\
GraphReconstructionBridge (TENSOR\_\allowbreak{}*, LASSERT) & \texttt{coq/\allowbreak{}kami\_\allowbreak{}hw/\allowbreak{}GraphReconstructionBridge.v} \\
ThieleCPU Kami specification & \texttt{coq/\allowbreak{}kami\_\allowbreak{}hw/\allowbreak{}ThieleCPUCore.v} \\
OCaml extraction observable surface & \texttt{coq/\allowbreak{}kernel/\allowbreak{}hardware\_\allowbreak{}bridge/\allowbreak{}OCamlExtractionBridge.v}: \texttt{ocaml\_\allowbreak{}bisimulation\_\allowbreak{}closure} \\
\end{longtable}
\endgroup

\chapter{Epistemological Classification}

The following table classifies every major result by its epistemological status.

\begin{remark}[How to read the tags]
This table is the document's principal credibility instrument, so four
classifications that a reader is most likely to trip over are stated
explicitly here.

\begin{itemize}
  \item \textbf{Tsirelson} carries two entries that are not the same result.
    The \emph{algebraic} bound is \textbf{(S)}: pure polynomial arithmetic,
    zero physics axioms. What is conditional is \emph{physical
    realizability}, the \texttt{A\_\allowbreak{}QM} bridge, which is its own line
    item and should never be folded into the algebra.
  \item \textbf{Born rule} likewise carries two entries for two results.
    \emph{Born Rule Uniqueness} is \textbf{(C)}, conditional on the
    linearity assumption. The \texttt{BornRule.v} relation is \textbf{(R)},
    because its $\mu$-consistency hypothesis goes unused in the proof.
  \item \textbf{Affine EFE} is \textbf{(R)}: reference-only geometry. The
    containing chapter is discrete \emph{Riemannian} geometry; the metric
    is positive-definite, so there is no causal structure to make a
    relativistic claim about (Remark~\ref{rem:not_spacetime}).
  \item \textbf{$\mu$-Chaitin} is \textbf{(C)}. It is a module-type
    contract with no instantiation anywhere in the corpus, and a contract
    with no instance is a conditional. Instantiating it is what would move it to \textbf{(S)}; the template is
    \texttt{Instance\_\allowbreak{}Kernel.v}, which does exactly this for NFI.
\end{itemize}
\end{remark}

\begingroup
\small
\begin{longtable}{@{}>{\raggedright\arraybackslash}p{0.30\textwidth}>{\raggedright\arraybackslash}p{0.10\textwidth}>{\raggedright\arraybackslash}p{\dimexpr0.60\textwidth-4\tabcolsep\relax}@{}}
\textbf{Result} & \textbf{Status} & \textbf{Key Assumption} \\
\hline
\multicolumn{3}{l}{\emph{Unconditional structural theorems about the machine}} \\
\hline
No Free Insight & \textbf{(S)} & Module type contract \\
$\mu$-Chaitin & \textbf{(C)} & Module-type contract with \emph{no instantiation anywhere in the corpus}; the chain is four assumed inequalities. Conditional until instantiated (cf.\ \texttt{Instance\_\allowbreak{}Kernel.v} for the NFI template) \\
Strict Subsumption & \textbf{(S)} & Partition structure \\
Oracle Cost (instantiation of universal NFI) & \textbf{(S)} & Definitional (soundness $:=$ cost $\ge 1$) \\
$\mu$-Conservation & \textbf{(S)} & Machine semantics only \\
$\mu$-Ledger Necessity & \textbf{(S)} & Machine semantics + explicit witnesses \\
$\mu$-Ledger Irredundancy (neither $\mu$ nor cert droppable) & \textbf{(S)} & Quantifies over all carriers/projections. \emph{Not} minimality: no order on projections is defined, so nothing is proved least (theorem \texttt{P\_\allowbreak{}full\_\allowbreak{}complete\_\allowbreak{}neither\_\allowbreak{}mu\_\allowbreak{}nor\_\allowbreak{}cert\_\allowbreak{}droppable}) \\
Receipt Integrity & \textbf{(S)} & Hash-chain model \\
Gauge / Noether & \textbf{(S)} & Definitional (record structure) \\
Three-Layer Observable Agreement (\emph{not} an isomorphism: no pair of mutually inverse maps is constructed; the content is transitivity of equality on a shared projection) & \textbf{(C)} & Backend decoders satisfy the shared observable projection contract \\
Discrete Angle-Defect Identity (planar triangulations) & \textbf{(S)} & Well-formed triangulated partition graph; equilateral model; not classical Gauss-Bonnet \\
Curvature from $\mu$-Gradients (uniform mass $\Rightarrow$ $G_{\mu\nu}=0$) & \textbf{(R)} & Reference only; $0=0$ in vacuum \\
Vacuum Implies Zero Stress-Energy & \textbf{(R)} & Reference only; $0=0$ in vacuum \\
Einstein Equations (Vacuum) & \textbf{(R)} & Reference only; not a claim about physics; $G=1/(8\pi)$ is a definition \\
Einstein Equations (General full tensor) & \textbf{(C)} & Named off-diagonal Ricci premise in \texttt{EinsteinEquationsFull.v} \\
Affine Off-Diagonal Ricci Zero (vertex 1) & \textbf{(R)} & Reference only; specific boundary 4-simplex configuration \\
Full Tensor EFE at Vertex 1 ($G=3T$) & \textbf{(R)} & Reference only; not a physical claim \\
Outer Vertex Characterization & \textbf{(R)} & Reference only \\
Discrete Differential Geometry & \textbf{(R)} & Definitions only; reference exploration \\
Partition Ops Free / Certification Non-Free & \textbf{(S)} & Machine semantics only \\
CHSH Trial Is Tier 0 & \textbf{(S)} & Machine semantics only \\
Three-Tier Observation Taxonomy & \textbf{(S)} & Machine semantics only \\
Locally Consistent $\Rightarrow$ Separable Coupling & \textbf{(S)} & Coupling definition \\
Local Coupling Bound ($|S| \le 2$) & \textbf{(S)} & Coupling definition \\
$S > 2$ Rules Out Local Couplings & \textbf{(S)} & Coupling definition \\
RTL Bisimulation (\texttt{driven\_\allowbreak{}step\_\allowbreak{}wf}) & \textbf{(S)} & Per-opcode Kami/\texttt{vm\_\allowbreak{}apply} agreement under \texttt{WFDrivenPrecondition}. (\texttt{rtl\_\allowbreak{}coverage\_\allowbreak{}partition}, $37{+}10{+}0=47$, is Peano arithmetic and is \emph{not} coverage evidence) \\
Slice-Coherent $\iff$ NPA Realizable & \textbf{(S)} & Polynomial arithmetic; standard-library assumptions \\
PR Box Not Slice-Coherent & \textbf{(S)} & Direct arithmetic on condition 1 \\
Elliptope Realizability (full correlator set) & \textbf{(S)} & Existential completion; stdlib reals only \\
Decidable Elliptope Gate Sound & \textbf{(S)} & Fraction-free $\mathbb{Z}$ arithmetic, two branches \\
Five Disciplines Are Unique Pointers & \textbf{(S)} & Closed under the global context; no axioms. Evidence for (PO-PUBLIC), not (PO-STRONG); see Section~\ref{sec:pointer_counterexamples} \\
Strong Pointer Criterion (PO-STRONG) Is False & \textbf{(S)} & Three machine-checked counterexamples (deniable authentication, symmetric MAC, object capabilities); forgery resistance does not imply metering \\
Tsirelson Is Thiele Boundary (forward + contrapositive) & \textbf{(S)} & From column contractivity; converse not claimed \\
Classical $P \subseteq \mu P$ & \textbf{(S)} & $\mu$-complexity definitions \\
\hline
\multicolumn{3}{l}{\emph{Conditional derivations (valid given stated axioms)}} \\
\hline
Born Rule Uniqueness & \textbf{(C)} & Linearity assumption (\texttt{coq/\allowbreak{}kernel/\allowbreak{}quantum/\allowbreak{}BornRule.v}, \texttt{BornRuleLinearity.v}) \\
Unitarity from $\mu=0$ & \textbf{(C)} & Info conservation (one direction only) \\
No-Cloning & \textbf{(C)} & Purity = information \\
Purification & \textbf{(C)} & Bloch sphere model \\
Landauer & \textbf{(C)} & Second law \emph{assumed} as record field \\
Self-Reference & \textbf{(C)} & Dimension = complexity \\

\hline
\multicolumn{3}{l}{\emph{Consistency relations (definitions verified, not predictions)}} \\
\hline
Information Causality & \textbf{(R)} & Circular: equivalence assumed as hypothesis \\
$\mu$-Geometry & \textbf{(R)} & Pseudometric (absolute value on $\mathbb{Z}$) \\
Spacetime Locality & \textbf{(R)} & Data isolation in partition graph \\
Born Rule (\texttt{BornRule.v}) & \textbf{(R)} & $\mu$-consistent hypothesis unused \\
\end{longtable}
\endgroup

\medskip
\noindent\textbf{Legend}: \textbf{(S)} = unconditional theorem about the machine; \textbf{(C)} = valid derivation conditional on stated axiom; \textbf{(R)} = internal consistency check.

\end{document}
